% Sumber LaTeX untuk buku ``Metode Aljabar'' dalam bahasa Tionghoa
% Hak cipta 2018  Li Wen-Wei.
% Izin diberikan untuk menyalin, menyebarluaskan, dan/atau mengubah
% dokumen ini berdasarkan ketentuan Creative Commons
% Atribusi 4.0 Internasional (CC BY 4.0)
% http://creativecommons.org/licenses/by/4.0/

% Untuk disertakan
\chapter{Dasar-Dasar Aljabar}\label{sec:algebras}
Banyak struktur gelanggang yang muncul dalam aljabar sekaligus merupakan ruang vektor atas suatu medan: penjumlahan gelanggang berasal dari penjumlahan ruang vektor, sedangkan perkalian $(x,y) \mapsto xy$ merupakan bentuk bilinear pada ruang vektor tersebut. Contoh khasnya ialah, atas medan $\Bbbk$, gelanggang matriks $n \times n$, yaitu $M_n(\Bbbk)$. Struktur semacam ini berulang kali dijumpai dalam praktik dan disebut ``aljabar'' atas medan. Secara umum, aljabar dapat dipandang sebagai gelanggang yang dilengkapi perkalian skalar oleh suatu gelanggang komutatif $R$, atau sebagai struktur gelanggang yang ditumpangkan pada sebuah $R$-modul; sudut pandang kedua mempunyai perumuman yang mendalam dalam kerangka kategori monoidal dan struktur kepang.

Selain konsep umum aljabar, termasuk operasi seperti hasil kali tensor dan perubahan medan, bab ini juga membahas:
\begin{compactenum}
	\item Integralitas dan keterhinggaan: konsep-konsep ini lazim dalam teori medan dan teori gelanggang komutatif, dan di sini akan ditangani secara umum. Salah satu hasil penting sekaligus menarik ialah Teorema Frobenius \ref{prop:division-R-algebra}, yang menyatakan bahwa, hingga isomorfisme, hanya ada tiga aljabar pembagian berdimensi hingga atas $\R$, yaitu $\R, \CC, \mathbb{H}$.
	\item Aljabar tensor beserta aljabar simetris dan aljabar eksterior yang diturunkan darinya: semuanya merupakan perangkat dasar geometri diferensial sekaligus kelanjutan alami aljabar linear. Misalnya, dengan aljabar eksterior, himpunan semua subruang berdimensi $k$ dari $\CC^n$ dapat diberi struktur geometri alami yang disebut varietas Grassmann. Untuk mengkajinya secara sistematis, konsep abstrak aljabar dan hasil kali tensor benar-benar diperlukan, dan konsep struktur bergradasi pun muncul secara alami. Bahasa kategori akan memperlihatkan kekuatannya yang khas dalam persoalan-persoalan ini.
	\item Teras dan norma juga merupakan konsep yang diperlukan dalam teori medan; pada hakikatnya keduanya hanyalah perumuman tertentu dari aljabar linear. Karena itu bab ini akan menatanya dalam kerangka umum dan membahas teori determinan atas gelanggang komutatif sebarang.
\end{compactenum}

\begin{wenxintishi}
	Karena pertimbangan penyusunan, bab ini tidak dapat memuat banyak contoh konkret yang menarik; pembaca sebaiknya terlebih dahulu menguasai sifat-sifat aljabar matriks. Akibatnya, gaya bab ini juga cukup formal; \S\ref{sec:integrality-finiteness} dan \S\ref{sec:Grassmannian} mungkin sedikit menyeimbangkannya. 

	Hasil kali tensor dan struktur bergradasi merupakan prasyarat untuk mempelajari aljabar tensor, dan untuk itu juga diperlukan bahasa kategori. Namun, pembaca dapat pula menempuh arah sebaliknya: karena aljabar tensor barangkali merupakan objek matematika yang lebih ``konkret'' (setidaknya lebih praktis), menjadikannya penghubung menuju semua prasyarat tadi juga merupakan pendekatan yang baik dan sesuai dengan urutan perkembangan historis. Hasil-hasil dalam \S\ref{sec:trace-norm-disc} hanya akan digunakan dalam kajian teori medan.
	
	Pembahasan terperinci tentang aljabar tensor selayaknya menjadi bagian dari mata kuliah aljabar linear; pembaca dapat merujuk pada buku teks baku seperti \cite[Bab 6]{Xi18}.
\end{wenxintishi}

\section{Aljabar atas Gelanggang Komutatif}\label{sec:algebra-def}
Selanjutnya, $R$ selalu menyatakan gelanggang komutatif tak nol. Modul atasnya tidak perlu dibedakan menjadi kiri dan kanan, meskipun notasi modul kiri akan dipakai. Semua aljabar dalam buku ini bersifat asosiatif dan mempunyai satuan.

Ingat bahwa setiap $R$-modul secara otomatis merupakan $(R,R)$-bimodul, cukup dengan mendefinisikan perkalian skalar dua sisi melalui $rmr' := rr'm$. Berikut ini kita akan memakai versi hasil kali seimbang untuk $(R,R)$-bimodul (Definisi \ref{def:balanced-product}), sebagaimana dijelaskan dalam Catatan \ref{rem:bimodule-balanced}.

\begin{definition}\label{def:algebra-naive}\index{daishu@aljabar (algebra)}
	Suatu \emph{aljabar} atas gelanggang $R$ ialah himpunan $A$ yang sekaligus mempunyai struktur gelanggang dan struktur $R$-modul, sedemikian sehingga penjumlahan gelanggang sama dengan penjumlahan $R$-modul dan perkalian gelanggang $(x,y) \mapsto xy$ merupakan hasil kali seimbang $A \times A \to A$ (dengan $A$ dipandang sebagai $(R,R)$-bimodul). Syarat terakhir ekuivalen dengan:
	\[ (rx)y = x(ry) = r(xy), \quad x,y \in A, \; r \in R. \]
	Homomorfisme antara aljabar $A_1$ dan $A_2$ berarti homomorfisme gelanggang $A_1 \to A_2$ yang $R$-linear; dengan cara serupa didefinisikan isomorfisme $R$-aljabar. Subaljabar ialah subgelanggang yang tertutup terhadap perkalian skalar, sedangkan ideal kiri/kanan dan ideal dua sisi didefinisikan secara serupa.\index{tongtai}\index{tonggou}
\end{definition}

\begin{convention}
	Agar teorinya seragam, kecuali dinyatakan lain kita tidak mengharuskan $R$-aljabar $A$ menjadi gelanggang tak nol. Manfaat konvensi ini mulai tampak dalam Definisi \ref{def:algebra-mod-diagram}.
\end{convention}

Semua kelas koset aditif dari aljabar $A$ modulo ideal dua sisi $I$ membentuk hasil bagi $A/I$. Dengan struktur gelanggang dan struktur modul hasil bagi yang telah tersedia, mudah dilihat bahwa $A/I$ merupakan aljabar, yang secara alami disebut aljabar hasil bagi\index{shang}. Berdasarkan hal ini, kita dapat membahas isomorfisme, hasil bagi, kernel homomorfisme, dan konsep-konsep lain untuk $R$-aljabar. Jika $A$ merupakan gelanggang pembagian sebagai gelanggang, maka $A$ disebut \emph{aljabar pembagian}. \index{daishu!aljabar pembagian (division algebra)}

Untuk ideal $I \subset R$, mudah dilihat bahwa $IA := \{ \sum_k t_k a_k : \forall k, \; t_k \in I, a_k \in A\}$ merupakan ideal dalam $A$.

Tuliskan $1$ untuk unsur satuan perkalian $A$, dan $Z_A$ untuk pusat $A$. Karena definisi menyiratkan
\[ rx = 1 \cdot rx = (r \cdot 1)x, \quad rx = rx \cdot 1 = x \cdot (r \cdot 1), \]
perkalian skalar dapat ditulis tanpa ambiguitas sebagai perkalian dari kiri maupun dari kanan. Selain itu, karena persamaan di atas memberikan $(rs) \cdot 1 = r(s \cdot 1) = (r \cdot 1) (s \cdot 1)$, pemetaan $r \mapsto r \cdot 1$ menghasilkan homomorfisme gelanggang $\sigma: R \to Z_A$. Sebaliknya, homomorfisme gelanggang $\sigma: R \to A$ memberikan struktur $R$-modul pada $A$ melalui $\sigma(r)a = ra$, dengan ruas kiri memakai perkalian gelanggang dan ruas kanan memakai perkalian skalar modul; struktur ini memenuhi Definisi \ref{def:algebra-naive} jika dan hanya jika $\Image(\sigma) \subset Z_A$.

\begin{proposition}\label{prop:algebra-as-homomorphism}
	Memberikan struktur $R$-aljabar pada gelanggang $A$ ekuivalen dengan memberikan homomorfisme gelanggang $\sigma: R \to Z_A$; korespondensinya ditentukan oleh $\sigma(r)a = ra$. Homomorfisme $\phi: A_1 \to A_2$ antara dua $R$-aljabar tidak lain adalah homomorfisme gelanggang $\phi: A_1 \to A_2$ yang membuat diagram berikut komutatif:
	\[\begin{tikzcd}[row sep=small]
		A_1 \arrow[rr, "\phi"] & & A_2 \\
		& R \arrow[lu, "\sigma_1"] \arrow[ru, "\sigma_2"'] &
	\end{tikzcd}\]
	dengan homomorfisme $\sigma_i: R \to Z_{A_i}$ menentukan struktur $R$-aljabar pada $A_i$ ($i=1,2$).
\end{proposition}
\begin{proof}
	Korespondensi antara aljabar dan homomorfisme $R \to Z_A$ telah dijelaskan di atas. Untuk homomorfisme gelanggang $\phi: A_1 \to A_2$ seperti di atas, tuliskan $1_{A_i}$ untuk unsur satuan $A_i$. Dari persamaan $\phi(r a) = \phi(r \cdot 1_{A_1}) \phi(a)$, terlihat bahwa $\phi$ bersifat $R$-linear jika dan hanya jika $\phi(r \cdot 1_{A_1}) = r \cdot 1_{A_2}$, dan inilah tepatnya syarat dalam pernyataan.
\end{proof}

Kategori semua $R$-aljabar dinotasikan dengan $R\dcate{Alg}$. Jika $R$ merupakan medan, kita dapat membicarakan dimensi suatu aljabar.
\begin{remark}
	Salah satu kasus khusus yang penting ialah $R = \Z$: $R$-aljabar tidak lain adalah gelanggang, atau dengan kata lain $\cate{Ring} \simeq \Z\dcate{Alg}$. Jadi, aljabar dapat dipandang sebagai perumuman gelanggang. Setidaknya ada dua penjelasan untuk hal ini.
	\begin{enumerate}[(i)]
		\item Grup abelian dan $\Z$-modul adalah hal yang sama; dengan kata lain, kategori $\cate{Ab}$ ekuivalen (bahkan isomorfik) dengan $\Z\dcate{Mod}$.
		\item Untuk setiap gelanggang $A$, terdapat tepat satu homomorfisme $\Z \to A$ seperti dalam \eqref{eqn:ring-struct-morphism}, dan citranya terletak dalam subgelanggang $Z_A$; dengan kata lain, $\Z$ adalah objek awal dalam kategori gelanggang $\cate{Ring}$.
	\end{enumerate}
\end{remark}

Contoh-contoh dasar aljabar meliputi:
\begin{itemize}
	\item Setiap gelanggang pembagian $D$ merupakan aljabar atas submedan primanya.
	\item Gelanggang monoid atas gelanggang komutatif $R$ (lihat \S\ref{sec:polynomial-ring}) merupakan $R$-aljabar; sebagai kasus khusus, gelanggang grup dan gelanggang polinomial yang dikenal tentu juga merupakan aljabar.
	\item Gelanggang matriks $n \times n$ atas gelanggang komutatif $R$, yaitu $M_n(R)$, merupakan $R$-aljabar.
\end{itemize}

\begin{example}\index{aljabar Calkin}
	Operator terbatas dari suatu ruang Hilbert terpisahkan berdimensi tak hingga $\mathcal{H}$ ke dirinya sendiri membentuk $\CC$-aljabar $B(\mathcal{H})$; himpunan operator kompak dinotasikan dengan $K(\mathcal{H})$. Dari teori operator dasar (misalnya \cite[\S 4.1]{Zh1}), diketahui bahwa $K(\mathcal{H})$ merupakan ideal dua sisi sejati dalam $B(\mathcal{H})$. Hasil bagi $B(\mathcal{H})/K(\mathcal{H})$ disebut aljabar Calkin, suatu objek alami dalam teori operator dan teori indeks; meninjau $B(\mathcal{H})/K(\mathcal{H})$ berarti meninjau kelas operator terbatas modulo perturbasi oleh operator kompak. Aljabar semacam ini juga membawa struktur tambahan yang berasal dari analisis (yang disebut aljabar $C^*$), dan menjadi salah satu contoh pertemuan antara aljabar dan analisis.
\end{example}

Pandang $R$ sebagai $R$-modul. Untuk setiap $R$-modul $M$,
\begin{inparaenum}[(i)]
	\item menentukan unsur $m$ dalam $M$ ekuivalen dengan menentukan homomorfisme modul $\eta: R \to M$, dengan hubungan $\eta(1) = m$;
	\item menentukan hasil kali seimbang $m: M \times M \to M$ dari $(R,R)$-bimodul ekuivalen dengan menentukan homomorfisme $R$-modul $\mu: M \dotimes{R} M \to M$ sedemikian sehingga $\mu(x \otimes y) = m(x,y)$. Ini merupakan terjemahan langsung dari sifat universal hasil kali tensor dalam Catatan \ref{rem:bimodule-balanced}. Karena $R$ komutatif, teori hasil kali tensor di sini jauh lebih sederhana daripada kasus dalam \S\ref{sec:module-tensor-prod}.
\end{inparaenum}

Sekarang Definisi \ref{def:algebra-naive} dapat ditulis ulang dalam bentuk diagram berikut. Peralihan antara kedua bentuk seharusnya mudah; pembaca dipersilakan memeriksanya. Untuk ringkasnya, kita tulis fungtor hasil kali tensor $R$-modul $- \dotimes{R} -$ sebagai $- \otimes -$.
\begin{definition}\label{def:algebra-mod-diagram}
	Aljabar atas gelanggang komutatif $R$ adalah struktur berikut:
	\begin{compactitem}
		\item suatu $R$-modul $A$,
		\item morfisme perkalian $\mu: A \otimes A \to A$,
		\item morfisme satuan atau unit $\eta: R \to A$,
	\end{compactitem}
	sedemikian sehingga semua diagram berikut komutatif:
	\begin{enumerate}[(i)]
		\item hukum asosiatif perkalian
		\[\begin{tikzcd}
		A \otimes (A \otimes A) \arrow[dash, rr, "\sim"] \arrow[d, "\identity \otimes \mu"'] & & (A \otimes A) \otimes A \arrow[d, "\mu \otimes \identity"] \\
		A \otimes A \arrow[r, "\mu"'] & A & A \otimes A \arrow[l, "\mu"] \\
		\end{tikzcd}\]
		dengan isomorfisme alami (``kendala asosiativitas'') $A \otimes (A \otimes A) \simeq (A \otimes A) \otimes A$ berasal dari Proposisi \ref{prop:tensor-assoc};
		\item hukum satuan kiri/kanan
		\[\begin{tikzcd}
		A \otimes R \arrow[r, "\identity \otimes \eta"] \arrow[rd, "\simeq"'] & A \otimes A \arrow[d, "\mu"] & R \otimes A \arrow[l, "\eta \otimes \identity"'] \arrow[ld, "\simeq"] \\
		& A &
		\end{tikzcd} \]
		dengan isomorfisme alami $R \otimes A \simeq A$ dan $A \otimes R \simeq A$ berasal dari Proposisi \ref{prop:tensor-unit}.
	\end{enumerate}
	Homomorfisme dari aljabar $(A_1, \mu_1, \eta_1)$ ke $(A_2, \mu_2, \eta_2)$ ialah homomorfisme modul $\phi: A_1 \to A_2$ yang membuat diagram-diagram berikut komutatif:
	\[\begin{tikzcd}
		A_1 \otimes A_1 \arrow[r, "\mu_1"] \arrow[d, "\phi \otimes \phi"'] & A_1 \arrow[d, "\phi"] \\
		A_2 \otimes A_2 \arrow[r, "\mu_2"'] & A_2
	\end{tikzcd}\quad\begin{tikzcd}
		A_1 \arrow[rr, "\phi"] & & A_2 \\
		& R \arrow[lu, "\eta_1"] \arrow[ru, "\eta_2"'] &
	\end{tikzcd}\]
\end{definition}
Subaljabar, ideal kiri dan kanan, serta konsep serupa dapat didefinisikan dengan cara yang sama memakai diagram dan submodul. Ciri khas pendekatan ini ialah bahwa, setelah kategori $R$-modul, fungtor hasil kali tensor $- \otimes -$, dan objek $R$ yang berperan sebagai ``satuan'' dianggap telah diberikan, definisi $R$-aljabar dapat dinyatakan sepenuhnya melalui panah dan komposisinya. Jika $A$ diambil sebagai modul nol, diperoleh aljabar nol.

\begin{example}
	Ingat Catatan \ref{rem:module-internal-Hom}: himpunan $\Hom$ antara $R$-modul mempunyai struktur $R$-modul alami, dan komposisi homomorfisme merupakan bentuk $R$-bilinear. Karena itu, semua endomorfisme suatu $R$-modul $M$ membentuk $R$-aljabar $\End_R(M)$.
\end{example}

\begin{example}\label{eg:matrix-algebra-1}\index{daishu!matriks}
	Dalam contoh sebelumnya, ambil modul bebas $M = R^{\oplus n}$. Aljabar endomorfismenya tidak lain adalah aljabar matriks $M_n(R)$ atas $R$. Teori matriks yang dikenal menyatakan bahwa $M_n(R)$ adalah $R$-modul bebas ber-rank $n^2$ dengan basis
	\[ E_{ij} := (a_{kl})_{1 \leq k,l \leq n}, \quad a_{kl} = \begin{cases} 1, & (k,l)=(i,j) \\ 0 & (k,l) \neq (i,j) \end{cases}. \]
	Karena perkalian aljabar bersifat bilinear, struktur $M_n(R)$ sepenuhnya ditentukan oleh perkalian unsur-unsur basis:
	\begin{equation}\label{eqn:E_ij}\begin{aligned}
		E_{ij} E_{kl} & = \begin{cases} E_{il}, & j=k, \\ 0, & j \neq k \end{cases} \\
		& = \delta_{j,k} E_{il};
	\end{aligned}\end{equation}
	di sini $\delta_{j,k}$ ialah notasi $\delta$ Kronecker: jika $j=k$, maka $\delta_{j,k}=1$; jika tidak, nilainya nol.
\end{example}

Secara umum, jika aljabar $A$ sebagai $R$-modul mempunyai basis $(e_i)_{i \in I}$ (basis selalu ada jika $R$ merupakan medan), maka struktur $A$ sepenuhnya ditentukan oleh perkalian basis
\[ e_i e_j = \sum_{k \in I} a^k_{i,j} e_k, \quad i,j \in I  \]
Data $(a^k_{i,j} \in R)_{i,j,k \in I}$ disebut \emph{konstanta struktur} $A$ terhadap basis tersebut.

Lebih lanjut, untuk setiap $R$-aljabar $A$ kita dapat meninjau gelanggang matriks $n \times n$, yaitu $M_n(A)$. Melalui $a \mapsto \biggl( \begin{smallmatrix} a & & \\ & \ddots & \\ & & a \end{smallmatrix} \biggr)$, $A$ tertanam sebagai subgelanggang $M_n(A)$, dan operasi matriks baku memberikan $Z_A = Z_{M_n(A)}$; khususnya, $M_n(A)$ juga merupakan $R$-aljabar. Dengan matriks $E_{ij}$ tetap didefinisikan seperti di atas, $M_n(A)$ adalah modul kiri bebas atas $A$:
\[ M_n(A) = \bigoplus_{1 \leq i,j \leq n} A E_{ij}, \]
Struktur $R$-aljabar pada $M_n(A)$ sepenuhnya ditentukan oleh dekomposisi di atas, struktur $R$-aljabar pada $A$, persamaan \eqref{eqn:E_ij}, serta $a E_{ij} = E_{ij} a$, dengan $a \in A$ dan $1 \leq i,j \leq n$.

Menurut Proposisi \ref{prop:End-matrix}, gelanggang $M_n(A)$ juga dapat dipahami sebagai $\End((A_A)^{\oplus n})$, dengan $A_A$ menyatakan $A$ sebagai modul kanan atas $A$ dan $\End = \End_{\cated{Mod}A}$. Ini memberi penafsiran intrinsik bagi $M_n(A)$ yang tidak bergantung pada pemilihan basis: ia adalah gelanggang endomorfisme dari modul kanan bebas ber-rank $n$ atas $A$. Karena perkalian skalar oleh $R \hookrightarrow Z_A$ pada modul kanan atas $A$ memberikan endomorfisme, struktur $R$-aljabar pada $M_n(A)$ juga menjadi jelas. Aljabar berbentuk $M_n(A)$ beserta subaljabarnya merupakan sumber gagasan yang tiada habisnya dalam teori gelanggang tak komutatif. Seperti dalam aljabar linear, sudut pandang intrinsik dan sudut pandang matriks harus dikuasai bersama.

\section{Integralitas, Keterhinggaan, dan Teorema Frobenius}\label{sec:integrality-finiteness}
Misalkan $A$ aljabar atas gelanggang komutatif $R$; dalam bagian ini selalu diasumsikan $A \neq \{0\}$. Untuk setiap $x \in A$, definisikan $R[x]$ sebagai subaljabar terkecil yang memuat $x$, atau subaljabar yang dibangkitkan oleh $x$. Evaluasi polinomial satu peubah dalam \eqref{eqn:polynomial-ev} memberikan homomorfisme surjektif \index[sym1]{$R[x]$}
\begin{align*}
	\text{ev}_x: R[X] & \longrightarrow R[x] \subset A \\
	\left( f(X) = \sum_k a_k X^k \right) & \longmapsto f(x) = \sum_k a_k x^k 
\end{align*}
Notasi $R[x]$ dengan demikian mempunyai arti yang intuitif. Secara lebih umum, untuk setiap $x,y, \ldots \in A$ dapat ditinjau subaljabar $R[x,y, \ldots]$ yang dibangkitkan oleh unsur-unsur tersebut. Biasanya kita hanya meninjau kasus ketika unsur-unsur itu saling komutatif; evaluasi polinomial kemudian tetap memberikan homomorfisme surjektif
\[ \text{ev}_{(x,y, \ldots)}: R[X, Y, \ldots] \twoheadrightarrow R[x,y,\ldots]. \]
Dengan demikian, $R[x,y, \ldots]$ selalu komutatif.

\begin{definition}[Integralitas]\label{def:integrality}\index{zhengyuan@unsur integral (integral element)}
	Jika untuk $x \in A$ terdapat $n \geq 1$ dan $a_0, \ldots, a_{n-1} \in R$ sedemikian sehingga
	\begin{equation}\label{eqn:integrality}
		x^n + a_{n-1} x^{n-1} + \cdots + a_0 = 0
	\end{equation}
	maka $x$ disebut integral atas $R$.
\end{definition}
Sebagai contoh, jika $R=\Z$ dan $A \subset \CC$, maka $x \in A$ integral jika dan hanya jika ia merupakan akar suatu polinomial monik berkoefisien bilangan bulat, yakni jika $x$ adalah bilangan bulat aljabar.

Suatu modul kiri $N$ atas gelanggang sebarang $S$ disebut setia jika $\text{ann}(N) = \{0\}$; dengan kata lain, $s \in S$ memenuhi $sN=\{0\}$ jika dan hanya jika $s=0$. Definisi yang sama berlaku bagi modul kanan. Melalui perkalian kiri, aljabar $A$ menjadi $R[x]$-modul setia; bahkan setiap submodul $A$ yang memuat $1$ juga setia.

\begin{theorem}\label{prop:integrality-finiteness}
	Untuk setiap $x \in A$, pernyataan-pernyataan berikut ekuivalen:
	\begin{enumerate}[(i)]
		\item $x$ integral atas $R$;
		\item $R[x]$ merupakan $R$-modul yang dibangkitkan secara hingga;
		\item $x$ termuat dalam suatu $R[x]$-submodul setia $M$ dari $A$, dengan $M$ dibangkitkan secara hingga sebagai $R$-modul.
	\end{enumerate}
	Sebagai kasus khusus dari (iii), jika $A$ dibangkitkan secara hingga sebagai $R$-modul, maka setiap $x \in A$ integral.
\end{theorem}
Pembuktian akan memakai teori determinan dan matriks adjoin atas gelanggang komutatif sebarang. Pembaca semestinya mengenal kasus medan, misalnya \cite[Proposisi 4.11]{Xi16}; kasus umum dapat dibuktikan dengan cara serupa dan akan dirinci dalam Teorema \ref{prop:matrix-det}.
\begin{proof}
	(i) $\implies$ (ii). Misalkan $x$ integral atas $R$. Dengan memakai \eqref{eqn:integrality} berulang kali, setiap polinomial dalam $x$ berkoefisien $R$ dapat ditulis ulang dalam bentuk berderajat $< n$. Jadi $R[x] = \sum_{k=0}^{n-1} R x^k$ dibangkitkan secara hingga.
	
	(ii) $\implies$ (iii) jelas.
	
	(iii) $\implies$ (i). Misalkan $x \in M = \sum_{k=1}^m R b_k$. Dari $xM \subset M$, terdapat matriks $T = (t_{ij})_{1 \leq i,j \leq m} \in M_m(R)$ sedemikian sehingga
	\[ x b_i = \sum_{j=1}^m t_{ij} b_j , \quad i = 1, \ldots, m \]
	yaitu persamaan matriks
	\begin{equation}\label{eqn:integrality-matrix}
		( x \cdot 1_{m \times m} - T) \begin{pmatrix} b_1 \\ \vdots \\ b_m \end{pmatrix}  = \begin{pmatrix} 0 \\ \vdots \\ 0 \end{pmatrix}.
	\end{equation}
	Untuk sementara, perkenalkan peubah formal $X$ dan tetapkan $S(X) := X \cdot 1_{m \times m} - T \in M_m(R[X])$. Tuliskan $S^\vee(X)$ untuk matriks adjoinnya. Maka $S^\vee(X) S(X) = P(X) \cdot 1_{m \times m}$, dengan $P(X) := \det(X \cdot 1_{m \times m} - T) \in R[X]$ polinomial karakteristik $T$, yang tentu monik berderajat $m$. Kalikan kedua ruas \eqref{eqn:integrality-matrix} dari kiri dengan $S^\vee(x) = \text{ev}_x(S^\vee(X))$. Dari persamaan $S^\vee(x)S(x) = P(x) \cdot 1_{m \times m}$ dalam $M_m(R[x])$, diperoleh
	\[ P(x) b_i = 0, \quad i=1, \ldots, m. \]
	Jadi $P(x) M=0$, dan sifat setia menyiratkan $P(x)=0$, yakni membuktikan integralitas.
\end{proof}

\begin{corollary}\label{prop:integrality-compositum}
	Misalkan $\{x_i\}_{i \in I}$ suatu keluarga unsur integral dalam $A$ yang saling komutatif terhadap perkalian. Maka setiap unsur dalam subaljabar $R[\{x_i\}_{i \in I}]$ yang mereka bangkitkan adalah integral. Sebagai akibatnya, jika $x, y \in A$ merupakan unsur integral yang saling komutatif dan $r,s \in R$, maka $rx + sy$ dan $xy$ juga integral.
\end{corollary}
\begin{proof}
	Karena setiap unsur tertentu dalam $R[\{x_i\}_{i \in I}]$ dapat dinyatakan dengan hanya sejumlah hingga $x_i$, cukup meninjau kasus sejumlah hingga unsur $x_1, \ldots, x_n$. Tinjau $R$-aljabar komutatif $R_m := R[x_1, \ldots, x_m] \subset A$, dengan $R_0 := R$. Untuk setiap $0 \leq m < n$:
	\begin{compactitem}
		\item $R_m \subset R_{m+1} = R_m[x_{m+1}]$,
		\item $x_{m+1}$ integral atas $R$, dan karena itu tentu integral atas $R_m$; maka (ii) dalam Teorema \ref{prop:integrality-finiteness} menyatakan bahwa $R_{m+1}$ dibangkitkan secara hingga sebagai $R_m$-modul.
	\end{compactitem}
	Dengan induksi, diperoleh bahwa $R_n$ dibangkitkan secara hingga sebagai $R$-modul. Bagian (iii) Teorema \ref{prop:integrality-finiteness} kemudian menyiratkan bahwa semua unsur $R_n$ integral (ambil $M = R_n$). Terbukti.
\end{proof}
\begin{remark}
	Untuk unsur integral $x$ dan $y$ beserta persamaan polinomial \eqref{eqn:integrality} yang dipenuhinya, pembuktian di atas hanya menunjukkan secara abstrak bahwa $x + y$ dan $xy$ integral. Jika ingin membangun secara eksplisit polinomial monik yang berakar $x+y$ atau $xy$, cara klasiknya ialah memakai \emph{resultan}.
\end{remark}

Dalam penerapan, kasus ketika $R$ merupakan medan sangat penting. Karena itu, berikut ini tetapkan suatu medan $\Bbbk$ dan biarkan $A$ menjadi $\Bbbk$-aljabar. Gelanggang $\Bbbk[X]$ adalah gelanggang ideal utama (Contoh \ref{eg:polynomial-PID}). Karena satu-satunya ideal sejati suatu medan ialah $\{0\}$, homomorfisme $\Bbbk \to Z_A$ pasti injektif, sehingga $\Bbbk$ dapat diidentifikasi dengan subgelanggang $A$.

\begin{definition}\label{def:algebraic-element}\index{daishuyuan@unsur aljabar (algebraic element)}
	Untuk unsur $x$ dalam $\Bbbk$-aljabar $A$, tinjau homomorfisme evaluasi $\text{ev}_x: \Bbbk[X] \to A$. Kernelnya $\Ker(\text{ev}_x)$ dapat ditulis sebagai ideal utama $(P_x)$.
	\begin{compactitem}
		\item Jika $P_x = 0$, maka $x$ bukan akar dalam $A$ dari polinomial tak nol mana pun; kita mengatakan bahwa atas $\Bbbk$, $x$ bersifat \emph{transenden};
		\item Jika $P_x \neq 0$, maka atas $\Bbbk$, $x$ disebut \emph{aljabar}. Setelah membagi $P_x$ dengan koefisien utamanya, kita dapat mengambil $P_x$ monik agar unik. Karena $P_x$ adalah polinomial monik berderajat terendah dalam $\Ker(\text{ev}_x)$, wajar menyebutnya \emph{polinomial minimal} dari $x$. \index{jixiaoduoxiangshi@polinomial minimal (minimal polynomial)} \index[sym1]{$P_x$}
	\end{compactitem}
\end{definition}
Pembahasan di atas juga menunjukkan bahwa atas medan $\Bbbk$, integralitas ekuivalen dengan sifat aljabar. Polinomial $P_x$ secara intrinsik mencirikan struktur subaljabar $\Bbbk[x] \simeq \Bbbk[X]/(P_x)$.

\begin{lemma}
	Misalkan $A$ suatu $\Bbbk$-aljabar. Untuk unsur aljabar $x$ dalam $A$,
	\[ x\; \text{invertibel dari kiri} \iff x\; \text{invertibel dalam $\Bbbk[x]$} \iff x\; \text{invertibel dari kanan}. \]
\end{lemma}
\begin{proof}
	Jelas bahwa $x \in \Bbbk[x]^\times$ menyiratkan $x$ invertibel dari kiri dan kanan dalam $A$. Sebaliknya, tuliskan polinomial minimal sebagai $P_x(X) = X^n + \cdots + a_1 X + a_0$. Jika $x$ invertibel dari kiri atau dari kanan, maka $a_0 \neq 0$; sebab jika tidak,
	\[ x(x^{n-1} + \cdots + a_1) = (x^{n-1} + \cdots + a_1)x = 0 \implies x^{n-1} + \cdots + a_1 = 0, \]
	maka $x$ menjadi akar suatu polinomial berderajat lebih rendah. Oleh karena itu, dari
	\[ x (x^{n-1} + \cdots + a_1) = -a_0 \in \Bbbk^\times \]
	diperoleh $x \in \Bbbk[x]^\times$.
\end{proof}

\begin{lemma}\label{prop:minimal-polynomial}
	Jika $\Bbbk$-aljabar $A$ tidak mempunyai pembagi nol selain unsur nol, maka $\Bbbk[x]$ merupakan medan untuk setiap unsur aljabar $x \in A$; selain itu, $P_x$ tak teruraikan dan $\dim_\Bbbk \Bbbk[x] = \deg P_x$.
\end{lemma}
\begin{proof}
	Tuliskan $\Bbbk$-subaljabar $\Bbbk[x]$ dari $A$ sebagai $\Bbbk[X]/(P_x)$. Karena gelanggang ini tidak mempunyai pembagi nol tak nol, $(P_x)$ merupakan ideal prima dalam $\Bbbk[X]$. Karena $\Bbbk[X]$ adalah gelanggang ideal utama, Teorema \ref{prop:PID-UFD} menyiratkan bahwa $P_x$ tak teruraikan dan $\Bbbk[X]/(P_x)$ merupakan medan. Sebagai ruang vektor atas $\Bbbk$, $\Bbbk[X]/(P_x)$ mempunyai basis jelas $\{X^i: 0 \leq i < \deg P_x \}$, sehingga dimensinya $\deg P_x$.
\end{proof}

Dalam $\Bbbk$-aljabar berdimensi hingga $A$, setiap unsur $x$ bersifat aljabar. Jika dipilih basis $a_1, \ldots, a_n$ dari $A$, terdapat pembenaman $\Bbbk$-aljabar
\begin{equation*}\begin{aligned}
	m: A & \longrightarrow \End_\Bbbk(A) \simeq M_n(\Bbbk) \\
	a & \mapsto \left[m_a: x \mapsto ax \right].
\end{aligned}\end{equation*}
Dengan demikian, $A$ dapat direalisasikan secara konkret di dalam aljabar matriks, sehingga beberapa sifat dapat direduksi ke kasus matriks, misalnya $xy=1 \iff yx=1$. Namun, keadaan semacam ini jarang dijumpai.

Aljabar berdimensi hingga atas $\Bbbk=\R$ adalah yang mula-mula menarik perhatian matematikawan. Contoh klasiknya ialah $\R$, $\CC$, $M_n(\R)$, dan, dari Contoh \ref{eg:Hamilton-quaternion}, aljabar kuaternion $\mathbb{H}$. Medan bilangan real $\R$ dan kompleks $\CC$ sudah dikenal; aljabar matriks $M_n(\R)$ tidak komutatif dan bukan aljabar pembagian, sedangkan sifat $\mathbb{H}$ sangat mencolok: ia adalah aljabar pembagian tak komutatif berdimensi 4, hasil luar biasa pertama dalam upaya memperumum medan bilangan kompleks. Pertanyaan alaminya ialah: hingga isomorfisme, adakah aljabar pembagian berdimensi hingga lain atas $\R$? Jawabannya tidak.

\begin{lemma}\label{prop:real-extension-C}
	Misalkan gelanggang $A$ tidak mempunyai pembagi nol selain unsur nol, dan $x \in A$.
	\begin{compactenum}[(i)]
		\item Jika $A$ merupakan $\CC$-aljabar dan $x$ aljabar atas $\CC$, maka $\CC[x]=\CC$;
		\item Jika $A$ merupakan $\R$-aljabar dan $x$ aljabar atas $\R$, maka $\R[x] = \R$ atau $\R[x] \simeq \CC$.
	\end{compactenum}
	Khususnya, jika semua unsur suatu $\CC$-aljabar $A$ bersifat aljabar, maka $A=\CC$.
\end{lemma}
\begin{proof}
	Tinjau polinomial minimal $P_x$ dari $x$ atas $\R$ atau $\CC$. Menurut Lemma \ref{prop:minimal-polynomial}, $P_x$ tak teruraikan. Dalam kasus $\CC$-aljabar, teorema dasar aljabar menyatakan bahwa polinomial kompleks tak teruraikan harus berderajat satu, sehingga $\CC[x]=\CC$. Dalam kasus $\R$-aljabar, jika $\deg P_x = 1$, maka $x \in \R$ dan $\R[x] = \R$. Jika tidak, $P_x$ adalah polinomial kuadrat tanpa akar real, $X^2 -2bX + c$; ini juga merupakan akibat teorema dasar aljabar. Berikut kita buktikan $\R[X]/(P_x) \simeq \CC$. Dengan melengkapkan kuadrat,
	\[ P_x(X) = (X-b)^2 - b^2 + c = (c - b^2)(Y^2 + 1), \quad Y := \frac{X-b}{\sqrt{c - b^2}}. \]
	Setiap $f(X) \in \R[X]$ dapat ditulis ulang sebagai polinomial dalam peubah $Y$, yang kita notasikan dengan $\tilde{f}(Y) \in \R[Y]$. Isomorfisme yang dicari diberikan oleh
	\[\begin{tikzcd}[row sep=tiny]
		\R[X]/(P_x) \arrow[r, "\sim"] & \R[Y]/(Y^2+1) = \R \oplus \R Y \arrow[r, "\sim"] & \CC \\
		f(X) \bmod (P_x) \arrow[mapsto, r] & \tilde{f}(Y) \bmod (Y^2+1) \arrow[mapsto, r] & \tilde{f}(i) \\
		& r + sY \bmod (Y^2 + 1) \arrow[mapsto, r] & r + si
	\end{tikzcd}\]
	dengan $i \in \CC$ suatu akar kuadrat dari $-1$. Pembaca dapat membandingkannya dengan Contoh \ref{eg:Cauchy-C}.
\end{proof}

\begin{theorem}[F.\ G.\ Frobenius, 1877]\label{prop:division-R-algebra}\index{Teorema Frobenius}
	Misalkan $D$ aljabar pembagian atas $\R$ dan setiap unsurnya bersifat aljabar. Maka $D$ isomorfik dengan $\R$, $\CC$, atau $\mathbb{H}$.
\end{theorem}
Perhatikan bahwa semua unsur aljabar berdimensi hingga bersifat aljabar. Pembuktian berikut diambil dari \cite[(13.12)]{Lam01}; pembuktian ini sekaligus menjelaskan mengapa definisi perkalian pada $\mathbb{H}$ sebenarnya muncul secara alami.
\begin{proof}[R.\ Palais]
	Kita boleh mengasumsikan $\dim_{\R}(D) \geq 2$ (jika tidak, $D=\R$). Lemma \ref{prop:real-extension-C} menunjukkan bahwa untuk setiap $x \in D \smallsetminus \R \implies \R[x] \simeq \CC$, sehingga sekurang-kurangnya satu salinan $\CC$ dapat ditanamkan ke dalam $D$. Tetapkan sebuah pembenaman $\CC \hookrightarrow D$ dan berikan pada $D$ struktur ruang vektor $\CC$ melalui perkalian kiri; tuliskan $i$ untuk satuan imajiner $\sqrt{-1}$. Definisikan
	\[ D^\pm := \left\{ x \in D : xi = \pm ix \right\} = \left\{ x \in D: ixi^{-1} = \pm x \right\}; \]
	Perhatikan bahwa $\Ad(i): x \mapsto ixi^{-1}$ sekaligus merupakan automorfisme gelanggang dan pemetaan $\CC$-linear, serta $\Ad(i)^2 = \identity_D$. Maka $D$ terurai menjadi ruang-ruang eigen untuk $\pm 1$:
	\begin{align*}
		D & = D^+ \oplus D^- \\
		x & = x^+ + x^-, \quad x^\pm = \frac{x \pm \Ad(i)(x)}{2}.
	\end{align*}
	Pertama-tama kita buktikan $D^+ = \CC$. Dari definisi, mudah diperiksa bahwa $D^+ \supset \CC$ dan $D^+$ merupakan $\CC$-aljabar. Karena semua unsurnya aljabar atas $\R$, unsur-unsur itu juga aljabar atas $\CC$; Lemma \ref{prop:real-extension-C} menyiratkan $D^+ = \CC$.

	Jika $D^- = \{0\}$, maka $D = \CC$. Jika tidak, ambil $j \in D^- \smallsetminus \{0\}$. Pemetaan perkalian kanan $t \mapsto tj$ memberikan homomorfisme injektif $\CC$-linear $D^- \to D^+$. Jadi $1 \leq \dim_{\CC} D^- \leq \dim_{\CC} D^+ = 1$, sehingga
	\begin{gather*}
		D^- = D^+ j, \\
		D = \R \oplus \R i \oplus \R j \oplus \R ij.
	\end{gather*}
	Karena polinomial minimal $j$ atas $\R$ berderajat dua dan tak teruraikan, $j^2 \in \R \oplus \R j$ serta $j^2 \notin \R_{>0}$. Di sisi lain, $j^2 \in D^+ \implies j^2 \in \CC = \R \oplus \R i$. Maka $j^2 \in \R_{<0}$. Setelah menskalakan $j$ dengan unsur yang sesuai dari $\R^\times$, kita boleh mengasumsikan $j^2 = -1$. Dengan demikian, perkalian pada $D$ sepenuhnya ditentukan oleh
	\[ i^2 = j^2 = -1, \quad ji=-ij. \]
	Membandingkan dengan Contoh \ref{eg:Hamilton-quaternion}, diperoleh $D \simeq \mathbb{H}$.
\end{proof}
Jika hukum asosiatif perkalian aljabar dilonggarkan lebih jauh, $\mathbb{H}$ dapat ditanamkan dalam suatu struktur yang disebut \emph{aljabar oktonion}, yaitu $\mathbb{O}$. Bertolak dari $\R$, salah satu cara untuk membangun $\CC$, $\mathbb{H}$, lalu $\mathbb{O}$ secara bertahap ialah proses Cayley--Dickson, juga disebut ``penggandaan''; latihan akan membahasnya lebih mendalam. \index{bayuanshu@oktonion (octonion)}

Teori aljabar pembagian berdimensi hingga atas medan yang lebih umum, seperti $\Q$, jauh lebih rumit dan berkaitan erat dengan berbagai persoalan dalam teori bilangan aljabar dan geometri. Kita akan kembali ke topik ini kelak.

\section{Hasil Kali Tensor Aljabar}\label{sec:algebra-tensor-product}
Untuk $R$-aljabar $A, B$ yang diberikan, bagian ini menjelaskan cara memberikan struktur $R$-aljabar baku pada $A \otimes B$ (meskipun bukan satu-satunya pilihan). Berbeda dari teori modul, di sini dapat ditinjau hasil kali tensor dari tak hingga banyak aljabar (Definisi--Teorema \ref{def:inf-tensor-product}). \index{zhangliangji!aljabar}

\begin{definition}
	Misalkan $A$ dan $B$ merupakan $R$-aljabar, dengan pemetaan perkalian masing-masing dinotasikan $\mu_A$, $\mu_B$, dan seterusnya. Definisikan hasil kali tensor keduanya sebagai $A \otimes B$ yang dilengkapi perkalian
	\begin{multline*}
		\mu_{A \otimes B}: (A \otimes B) \otimes (A \otimes B) \rightiso (A \otimes (B \otimes A)) \otimes B \rightiso (A \otimes (A \otimes B)) \otimes B \\
		\rightiso (A \otimes A) \otimes (B \otimes B) \xrightarrow{\mu_A \otimes \mu_B} A \otimes B
	\end{multline*}
	dengan isomorfisme tanpa label seperti biasa merupakan kendala asosiativitas dan kendala komutativitas untuk $\otimes$ pada $R$-modul, serta dengan satuan
	\[ \eta_{A \otimes B}: R \rightiso R \otimes R \xrightarrow{\eta_A \otimes \eta_B} A \otimes B, \]
	dengan isomorfisme tanpa label berasal dari Proposisi \ref{prop:tensor-unit}, atau secara konkret diberikan oleh $r \mapsto r \otimes 1 = 1 \otimes r$. Jika operasi dalam $A \otimes B$ dijelaskan dengan keluarga pembangkit, rumusnya ialah
	\begin{gather*}
	(a \otimes b) \cdot (a' \otimes b') = aa' \otimes bb', \\
	1_{A \otimes B} = 1_A \otimes 1_B.
	\end{gather*}
\end{definition}
Sifat-sifat aljabar tersebut dapat diperiksa secara rutin meskipun agak panjang, dan rinciannya dihilangkan.

\begin{lemma}\label{prop:algebra-otimes-comm}
	Dalam situasi di atas, kendala komutativitas $c(A, B)$ memberikan isomorfisme alami $R$-aljabar $A \otimes B \rightiso B \otimes A$.
\end{lemma}
\begin{proof}
	Bahwa $c(A, B)$ merupakan isomorfisme aljabar dapat diperiksa melalui diagram dan sifat-sifat berbagai kendala (lihat \S\ref{sec:braiding}), atau langsung pada tingkat unsur sebagai berikut:
	\begin{gather*}
	(a \otimes b) \cdot (a' \otimes b') = aa' \otimes bb' \xrightarrow{c(A, B)} bb' \otimes aa' = c(A,B)(a \otimes b) \cdot c(A,B)(a' \otimes b'), \\
	c(A,B)(1_{A \otimes B}) = c(A,B)(1_A \otimes 1_B) = 1_B \otimes 1_A = 1_{B \otimes A}.
	\end{gather*}
\end{proof}

Perhatikan pula bahwa pembatasan $A \times B \to A \otimes B$ pada $A \times {1_B}$ dan ${1_A} \times B$ memberikan homomorfisme modul alami $\iota_A: A \to A \otimes B$ dan $\iota_B: B \to A \otimes B$. Citra keduanya dapat ditulis sebagai $A \otimes 1$ dan $1 \otimes B$, dan bersama-sama membangkitkan $A \otimes B$. Hasil berikut membenarkan notasi tersebut.
\begin{lemma}
	Jika $R$ merupakan medan, maka $\iota_A$ dan $\iota_B$ adalah pembenaman.
\end{lemma}
\begin{proof}
	Pilih basis $J_A$ dan $J_B$ masing-masing untuk $A$ dan $B$ atas $R$. Proposisi \ref{prop:tensor-direct-sum} memberikan isomorfisme alami ruang vektor atas $R$, $A \dotimes{R} B \simeq \bigoplus_{\substack{a \in J_A \\ b \in J_B }} Ra \dotimes{R} Rb$. Kita boleh mengasumsikan $1_A \in J_A$ dan $1_B \in J_B$. Dengan demikian, $\iota_A$ dan $\iota_B$ masing-masing menanamkan $A$ dan $B$ sebagai suku jumlah langsung.
\end{proof}

Dari sini diperoleh karakterisasi sifat universal $A \otimes B$.
\begin{proposition}\label{prop:algebra-otimes-univ}
	Dalam situasi di atas, $\iota_A$ dan $\iota_B$ merupakan homomorfisme aljabar, dan citra keduanya saling komutatif dalam $A \otimes B$. Data $(A \otimes B, \iota_A, \iota_B)$ memenuhi sifat universal berikut: untuk setiap data $(C, f_A, f_B)$, dengan $C$ suatu $R$-aljabar serta $f_A: A \to C$ dan $f_B: B \to C$ mempunyai citra yang saling komutatif, terdapat tepat satu homomorfisme aljabar $\phi: A \otimes B \to C$ yang membuat diagram berikut komutatif:
	\[\begin{tikzcd}
	A \arrow[r, "\iota_A"] \arrow[rd, "f_A"'] & A \otimes B \arrow[d, "\exists! \phi"] & B \arrow[l, "\iota_B"'] \arrow[ld, "f_B"] \\
	& C &
	\end{tikzcd}\]
	Hingga isomorfisme tunggal, sifat di atas mencirikan $A \otimes B$ secara unik.
\end{proposition}
\begin{proof}
	Mudah dilihat bahwa $\iota_A$ dan $\iota_B$ merupakan homomorfisme aljabar; dari $(a \otimes 1)(1 \otimes b) = a \otimes b = (1 \otimes b)(a \otimes 1)$ terlihat bahwa $A \otimes 1$ dan $1 \otimes B$ saling komutatif. Sekarang berikan data $(C, f_A, f_B)$. Diagram komutatif dalam pernyataan menunjukkan bahwa komposisi $A \times B \to A \otimes B \xrightarrow{\phi} C$ harus berupa $(a, b) \mapsto f_A(a)f_B(b)$. Pemetaan ini merupakan hasil kali seimbang, sehingga menentukan $\phi$ secara unik. Untuk keberadaan $\phi$, hasil kali seimbang $(a, b) \mapsto f_A(a)f_B(b)$ memberikan $\phi: A \otimes B \to C$, dan syarat bahwa citra $f_A$ dan $f_B$ saling komutatif tepat merupakan syarat yang diperlukan agar $\phi$ menjadi homomorfisme aljabar.
\end{proof}
Pembaca yang mempunyai cukup banyak kertas dapat mencoba menyatakan ulang sifat-sifat dan pembuktian ini dengan diagram.

\begin{remark}\label{rem:bimodule-as-module}\index{shuangmo}
	Sekarang kita dapat menjelaskan cara memasukkan bimodul ke dalam kerangka modul kiri atau modul kanan; cukup gunakan hasil kali tensor $\Z$-aljabar. Misalkan $S, T$ gelanggang, yakni $\Z$-aljabar. Maka terdapat isomorfisme kategori
	\[ \left(S \dotimes{\Z} T^\text{op}\right)\dcate{Mod} \simeq (S,T)\dcate{Mod} \simeq \cated{Mod}\left(S^\text{op} \dotimes{\Z} T\right). \]
	Hal ini karena memberikan struktur $(S,T)$-bimodul pada grup abelian $(M,+)$ ekuivalen dengan menentukan homomorfisme gelanggang $S \to \End(M)$ dan $T^\text{op} \to \End(M)$ yang citranya saling komutatif (bandingkan Catatan \ref{rem:module-multiplication}); ini sama dengan menentukan homomorfisme gelanggang $S \dotimes{\Z} T^\text{op} \to \End(M)$ atau $\left( S^\text{op} \dotimes{\Z} T \right)^\text{op} \to \End(M)$.
\end{remark}

\begin{remark}\label{rem:algebra-otimes-zero}
	Hasil kali tensor dua $R$-aljabar tak nol dapat menjadi aljabar nol. Misalnya, ambil $a, b \in \Z_{>1}$ yang saling prima. Terdapat bilangan bulat $u,v$ dengan $1 = ua+vb$, sehingga $\Z/a\Z \dotimes{\Z} \Z/b\Z = 0$. Namun, jika $A$ dan $B$ keduanya merupakan $R$-modul bebas tak nol, Korolari \ref{prop:module-tensor-free} memastikan bahwa $A \dotimes{R} B$ bebas dengan rank sama dengan hasil kali rank $A$ dan $B$, sehingga tak nol. Jika $R$ merupakan medan, asumsi ini selalu terpenuhi. Argumen tersebut langsung meluas ke hasil kali tensor tak hingga yang akan dibahas kemudian (Definisi--Teorema \ref{def:inf-tensor-product}).
\end{remark}

\begin{example}
	Kembali ke contoh dasar aljabar tak komutatif, yaitu aljabar matriks. Berikut ini tetapkan $\otimes := \otimes_R$. Kita akan membuktikan bahwa untuk setiap $R$-aljabar $A$ dan $B$ terdapat isomorfisme
	\begin{align}
		M_n(A) \otimes B & \rightiso M_n(A \otimes B) \label{eqn:matrix-tensor-1} \\
		M_n(R) \otimes M_m(R) & \rightiso M_{nm}(R). \label{eqn:matrix-tensor-2}
	\end{align}
	Dengan menggabungkan keduanya, diperoleh
	\begin{gather}\label{eqn:matrix-tensor-3}
		M_n(A) \otimes M_m(B) \simeq A \otimes M_n(R) \otimes M_m(R) \otimes B \simeq M_{nm}(A \otimes B).
	\end{gather}
	Pertama-tama tinjau \eqref{eqn:matrix-tensor-1}. Pemetaan tersebut tidak lain adalah $(a_{ij})_{i,j} \otimes b \mapsto (a_{ij} \otimes b)_{i,j}$. Proposisi \ref{prop:tensor-direct-sum}, Contoh \ref{eg:matrix-algebra-1}, dan pembahasan sesudahnya menunjukkan bahwa $M_n(A) \otimes B$ merupakan modul kiri bebas atas $A \otimes B$,
	\[ M_n(A) \otimes B = \bigoplus_{i,j} (A \otimes B) (E_{ij} \otimes 1), \]
	dan $M_n(A \otimes B)$ juga bebas: $M_n(A \otimes B) = \bigoplus_{i,j} (A \otimes B) E_{ij}$. Untuk setiap $i,j$, petakan $(a \otimes b) E_{ij}$ ke $(a \otimes b) (E_{ij} \otimes 1)$. Bijeksi ini jelas mempertahankan struktur perkalian, sehingga \eqref{eqn:matrix-tensor-1} memang suatu isomorfisme.
		
	Selanjutnya jelaskan \eqref{eqn:matrix-tensor-2}. Secara abstrak, misalkan $V$ dan $W$ masing-masing merupakan $R$-modul bebas ber-rank $n$ dan $m$. Proposisi \ref{prop:tensor-direct-sum} menyiratkan bahwa $V \otimes W$ adalah $R$-modul bebas ber-rank $nm$, sedangkan fungtorialitas hasil kali tensor memberikan homomorfisme $R$-aljabar $\End(V) \otimes \End(W) \to \End(V \otimes W)$.
	
	Pilih basis sehingga dapat diasumsikan $V = R^{\oplus n}$ dan $W = R^{\oplus m}$. Dengan demikian, $\End(V) \otimes \End(W)$ dapat diidentifikasi dengan $R$-modul bebas
	\[ M_n(R) \otimes M_m(R) = \bigoplus_{\substack{1 \leq i,j \leq n \\ 1 \leq r,s \leq m}} R E_{ij} \otimes E_{rs}  \]
	Sebagai $R$-aljabar, perkaliannya sepenuhnya ditentukan oleh $(E_{ij} \otimes E_{rs}) \cdot (E_{kl} \otimes E_{tu}) = \delta_{j,k} \delta_{s,t} E_{il} \otimes E_{ru}$. Pilih suatu bijeksi
	\[ f: \{1, \ldots, n\} \times \{1, \ldots, m\} \xrightarrow{1:1} \{1, \ldots, nm\} \]
	maka $V \otimes W$ dapat diidentifikasi dengan $R^{\oplus nm}$, dan mudah dilihat bahwa diagram
	\[\begin{tikzcd}
		\End(V) \otimes \End(W) \arrow[d, "\simeq"'] \arrow[r] & \End(V \otimes W) \arrow[d, "\simeq"] \\
		M_n(R) \otimes M_m(R) \arrow[r] &  M_{nm}(R) \\
		E_{ij} \otimes E_{rs} \arrow[mapsto, r] \arrow[phantom, sloped, u, "\in" description] & E_{f(i,r), f(j,s)} \arrow[phantom, sloped, u, "\in" description]
	\end{tikzcd}\]
	komutatif, sementara baris kedua jelas merupakan isomorfisme. Ini membangun isomorfisme \eqref{eqn:matrix-tensor-2} yang dicari.
\end{example}

Kembali ke pembahasan utama. Definisi dan sifat universal hasil kali tensor aljabar mempunyai perumuman langsung ke kasus banyak peubah $\bigotimes_{i \in I} A_i$, dengan $I$ himpunan hingga. Konstruksi ini jelas bersifat fungtorial. Misalnya, untuk keluarga homomorfisme $f_i: A_i \to A'_i$, terdapat pemetaan alami $\bigotimes_{i \in I} f_i : \bigotimes_{i \in I} A_i \to \bigotimes_{i \in I} A'_i$ yang membuat diagram
\[\begin{tikzcd}[row sep=small]
	A_j \arrow[r, "f_j"] \arrow[d] & A'_j \arrow[d] \\
	\bigotimes_i A_i \arrow[r, "{\bigotimes_i f_i}"'] & \bigotimes_i A'_i
\end{tikzcd}\]
komutatif untuk semua $j \in I$, dan demikian seterusnya.

Sifat universal dalam Proposisi \ref{prop:algebra-otimes-univ} dapat diperumum kepada sebarang keluarga $(f_i: A_i \to C)_{i \in I}$, dengan $I$ mungkin tak hingga. Untuk memberikan objek universal terkait, yakni hasil kali tensor tak hingga, perhatikan hal berikut. Jika $J \subset I$ adalah dua himpunan hingga, dengan meniru konstruksi $\iota_A$ sebelumnya kita memperoleh $f_{JI}: \bigotimes_{i \in J} A_i \to \bigotimes_{i \in I} A_i$, dan $K \subset J \subset I \implies f_{JI} f_{KJ} = f_{KI}$. Sekarang biarkan $I$ sebarang himpunan dan berikan keluarga $R$-aljabar $(A_i)_{i \in I}$. Perhatikan bahwa semua himpunan bagian hingga dari $I$, dengan relasi $\subset$, membentuk poset terarah ke atas sebagaimana didefinisikan dalam \S\ref{sec:group-limit}. Konstruksi $\varinjlim$ gelanggang atas poset terarah ke atas telah diberikan dalam Proposisi \ref{prop:ring-filtrant-limit}, dan tanpa kesulitan dapat diperluas ke kategori $R$-aljabar.

\begin{definition-theorem}\label{def:inf-tensor-product}
	Untuk setiap keluarga $R$-aljabar $(A_i)_{i \in I}$, definisikan
	\[ \bigotimes_{i \in I} A_i := \varinjlim_{\substack{J \subset I \\ \text{himpunan bagian hingga}}} \bigotimes_{i \in J} A_i \]
	dengan, untuk $K \subset J$, morfisme yang digunakan dalam pembentukan limit ialah $f_{KJ}: \bigotimes_{i \in K} A_i \to \bigotimes_{i \in J} A_i$. Objek ini dilengkapi keluarga homomorfisme $\iota_j: A_j \to \bigotimes_{i \in I} A_i$. Data $(\bigotimes_{i \in I} A_i, (\iota_i)_{i \in I})$ memenuhi sifat universal yang serupa dengan Proposisi \ref{prop:algebra-otimes-univ}.
\end{definition-theorem}
\begin{proof}[Sketsa]
	Mudah diperiksa bahwa $\iota_{i,J}: A_i \to \bigotimes_{i \in J} A_i$ kompatibel dengan keluarga homomorfisme $f_{KJ}$: $i \in K \subset J \implies f_{KJ} \circ \iota_{i,K} = \iota_{i,J}$. Karena itu kita memperoleh $\iota_i: A_i \to \bigotimes_{i \in I} A_i$; pemeriksaan bahwa semuanya homomorfisme aljabar bersifat rutin. Sifat universal juga direduksi ke kasus hasil kali tensor hingga: untuk data $(C, (f_i: A_i \to C)_{i \in I})$, pemetaan yang dicari $\phi: \bigotimes_i A_i \to C$ harus dan hanya dapat ditentukan oleh pemetaan $\phi_J: \bigotimes_{i \in J} A_i \to C$ melalui $\varinjlim_J$.
\end{proof}

\begin{corollary}\label{prop:tensor-alg-cocomplete}
	Hasil kali tensor memberikan koproduk dalam kategori $R\dcate{CAlg}$ dari $R$-aljabar komutatif. Kategori $R\dcate{CAlg}$ kokomplet.
\end{corollary}
\begin{proof}
	Mudah dibuktikan bahwa hasil kali tensor $R$-aljabar komutatif tetap komutatif; ini diserahkan sebagai latihan. Dalam sifat universal hasil kali tensor pada Proposisi \ref{prop:algebra-otimes-univ} (atau versi tak hingganya), jika hanya $R$-aljabar komutatif $C$ yang diperhitungkan, yang diperoleh tepat sifat universal koproduk. Di sisi lain, dalam $R\dcate{CAlg}$, setiap pasangan morfisme $f, g: A \to B$ mempunyai koekualiser $\Coker(f,g)$, yakni hasil bagi $B$ oleh ideal yang dibangkitkan oleh $\{f(a)-g(a): a \in A \}$ (aljabar nol diperbolehkan). Dari Teorema \ref{prop:limit-buildingblocks}, keberadaan $\varinjlim$ umum langsung diperoleh.
\end{proof}

Terakhir, kita memakai hasil kali tensor untuk membangun \emph{perubahan basis} aljabar\index{jibianhuan}. Ingat kembali konstruksi sederhana dalam \S\ref{sec:change-of-rings}: tetapkan homomorfisme gelanggang komutatif $\phi: R \to S$ dan pandang $S$ sebagai $R$-modul. Untuk setiap $R$-aljabar $A$, hasil kali tensor $A \dotimes{R} S$ secara alami mempunyai struktur $S$-aljabar. Menurut sudut pandang Proposisi \ref{prop:algebra-as-homomorphism}, hal ini karena citra $\iota_S: S \to A \dotimes{R} S$ terletak dalam pusat $A \dotimes{R} S$. % Semua konstruksi ini jelas fungtorial dalam peubah $A$.

Sebaliknya, menarik kembali perkalian skalar suatu $S$-aljabar $B$ melalui $\phi$ menghasilkan suatu $R$-aljabar. Jika hanya struktur modul yang diperhatikan, kedua operasi ini tepat merupakan fungtor dalam Korolari \ref{prop:IP-vs-F}, yaitu $\begin{tikzcd} R\dcate{Mod} \arrow[r, yshift=0.5ex, "{P_{R \to S}}"] & S\dcate{Mod} \arrow[l, yshift=-0.5ex, "{\mathcal{F}_{R \to S}}"] \end{tikzcd}$. Dengan memasukkan perkalian, pembahasan di atas memberikan sepasang fungtor $\begin{tikzcd} R\dcate{Alg} \arrow[r, yshift=0.5ex, "{P_{R \to S}}"] & S\dcate{Alg} \arrow[l, yshift=-0.5ex, "{\mathcal{F}_{R \to S}}"] \end{tikzcd}$.

\begin{proposition}\label{prop:IP-vs-F-alg}
	Fungtor-fungtor di atas membentuk pasangan adjoin $(P_{R \to S}, \mathcal{F}_{R \to S})$ antara kategori $R$-aljabar dan $S$-aljabar. Untuk homomorfisme $Q \to R \to S$, terdapat isomorfisme $P_{R \to S} \circ P_{Q \to R} \simeq P_{Q \to S}$ dan persamaan $\mathcal{F}_{Q \to R} \circ \mathcal{F}_{R \to S} = \mathcal{F}_{Q \to S}$.
\end{proposition}
\begin{proof}
	Misalkan $A$ dan $B$ masing-masing merupakan $R$-aljabar dan $S$-aljabar.	Kita mengingat kembali, untuk $(P_{R \to S}, \mathcal{F}_{R \to S})$, isomorfisme adjoin pada tingkat modul dalam \eqref{eqn:IP-adjunction-explicit-1}. Isomorfisme ini memetakan homomorfisme $R$-modul $f: A \to B$ ke homomorfisme $S$-modul $f': A \dotimes{R} S \to B$, dengan $f'(a \otimes s) = f(a)s \in B$; sebaliknya, komposisi $f'$ dengan $A \to A \dotimes{R} S$ mengembalikan $f$. Untuk membuktikan bahwa kedua fungtor itu adjoin pada tingkat aljabar, cukup dibuktikan bahwa korespondensi tadi memetakan homomorfisme aljabar ke homomorfisme aljabar, dan hal ini langsung.
	
	Sifat adjoin dan sifat-sifat yang berkaitan dengan komposisi homomorfisme $Q \to R \to S$ dapat ditangani dengan cara yang sama seperti Lemma \ref{prop:IP-transitivity}.
\end{proof}

\begin{proposition}\label{prop:algebra-base-change-monoidal}
	Fungtor $P_{R \to S}$ merupakan fungtor monoidal: terdapat isomorfisme alami $S$-aljabar
	\[ P_{R \to S}(A) \dotimes{S} P_{R \to S}(B) \rightiso P_{R \to S}(A \dotimes{R} B). \]
\end{proposition}
\begin{proof}
	Hal utamanya ialah menerapkan
	\begin{align*}
		(A \dotimes{R} S) \dotimes{S} (B \dotimes{R} S) & \longrightiso (A \dotimes{R} B) \dotimes{R} S \\
		(a \otimes s) \otimes (b \otimes t) & \longmapsto (a \otimes b) \otimes st, \\
		(a \otimes s) \otimes (b \otimes 1) & \longmapsfrom (a \otimes b) \otimes s
	\end{align*}
	Rinciannya agak panjang tetapi tidak sulit, dan diserahkan kepada pembaca.
\end{proof}

Suatu aljabar $A$ atas gelanggang $R$ dapat dipahami sebagai keluarga aljabar atas medan yang berubah secara ``aljabar'' dengan ruang parameter spektrum ideal maksimal $\MaxSpec(R)$. Untuk setiap ideal maksimal $\mathfrak{m}$, aljabar terkait atas medan $R/\mathfrak{m}$ ialah perubahan basis $P_{R \to R/\mathfrak{m}} (A) = A \dotimes{R} R/\mathfrak{m}$, yang juga disebut reduksi $A$ $\bmod\; \mathfrak{m}$. Istilah ini dibenarkan oleh hasil berikut.
\begin{lemma}
	Misalkan $A_1$ dan $A_2$ merupakan $R$-aljabar, sedangkan $\mathfrak{a}_i \subset A_i$ ideal. Maka terdapat isomorfisme alami
	\[ (A_1/\mathfrak{a}_1) \dotimes{R} (A_2/\mathfrak{a}_2) \rightiso A_1 \dotimes{R} A_2 \bigg/ \left( \mathfrak{a}_1 \dotimes{R} A_2 + A_1 \dotimes{R} \mathfrak{a}_2\right), \]
	Agar notasinya ringkas, di sini kita mengidentifikasi $\mathfrak{a}_1 \otimes A_2$ dan $A_1 \otimes \mathfrak{a}_2$ dengan citranya di dalam $A_1 \otimes A_2$.
	
	Sebagai akibatnya, untuk setiap ideal $I \subset R$ dan $R$-aljabar $A$, terdapat isomorfisme alami $A \dotimes{R} (R/I) \rightiso A/IA$.
\end{lemma}
\begin{proof}
	Untuk pernyataan pertama, tuliskan $\mathscr{A} := A_1 \otimes A_2 \big/ \left( \mathfrak{a}_1 \otimes A_2 + A_1 \otimes \mathfrak{a}_2\right)$. Bagi homomorfisme yang jelas $\iota_i: A_i/\mathfrak{a}_i \to \mathscr{A}$ ($i=1,2$), kita perlu memeriksa sifat universal dalam Proposisi \ref{prop:algebra-otimes-univ}; ini merupakan latihan langsung.

	Untuk pernyataan kedua, ambil $A_1 = A$, $\mathfrak{a}_1=\{0\}$, $A_2 = R$, dan $\mathfrak{a}_2 = I$, lalu gunakan isomorfisme $A \dotimes{R} R \rightiso A$, yang memetakan citra $A \otimes I$ ke $IA$.
\end{proof}

\section{Aljabar Bergradasi}
Struktur bergradasi pada aljabar muncul secara alami dalam penerapan; contoh awalnya ialah aljabar polinomial $R[X] = \bigoplus_{n \geq 0} R X^n$. Kita juga akan mempelajari hasil kali tensor antara aljabar bergradasi, sebagai perumuman konstruksi dalam \S\ref{sec:algebra-tensor-product}.

\begin{definition}[Modul dan aljabar bergradasi]\label{def:graded-module-alg}\index{mo!bergradasi (graded)}\index{daishu!bergradasi (graded)}\index{qiciyuan@unsur homogen (homogeneous element)}
	Misalkan $I$ monoid komutatif, dengan $+$ menyatakan operasi biner dan $0$ unsur satuan.
	\begin{itemize}
		\item Suatu $I$-\emph{modul bergradasi} atas gelanggang komutatif $R$ ialah modul $M$ yang dilengkapi dekomposisi jumlah langsung $M = \bigoplus_{i \in I} M_i$. Semua $I$-modul bergradasi membentuk kategori monoidal $(R\dcate{Mod}_I, \otimes)$: homomorfisme bergradasi dari $M$ ke $N$ berarti homomorfisme yang memenuhi $\forall i \; \varphi(M_i) \subset N_i$, sedangkan hasil kali tensor $M \otimes N$ membawa struktur bergradasi alami (ingat bahwa Proposisi \ref{prop:tensor-direct-sum} menyatakan bahwa hasil kali tensor mempertahankan jumlah langsung)
		\[ (M \otimes N)_k = \bigoplus_{\substack{i,j \in I \\ i+j=k}} M_i \otimes N_j. \]
		Unsur $x \in M_i \smallsetminus \{0\}$ disebut unsur \emph{homogen} berderajat $i$ dalam $M$, dan ditulis pula $\deg(x)=i$; derajat $0$ tidak didefinisikan. \emph{Submodul bergradasi} berarti submodul dari $M$ yang memenuhi $N = \bigoplus_i (N \cap M_i)$.
		\item Suatu $I$-\emph{aljabar bergradasi} atas gelanggang komutatif $R$ ialah $R$-aljabar yang dilengkapi dekomposisi jumlah langsung $A = \bigoplus_{i \in I} A_i$, dengan perkalian memenuhi $A_i \cdot A_j \subset A_{i+j}$ dan $1 \in A_0$; jadi $A_0$ adalah subaljabar. Dengan cara serupa, $I$-aljabar bergradasi membentuk kategori $R\dcate{Alg}_I$, dan homomorfisme $\varphi: A \to B$ harus sekaligus merupakan homomorfisme $I$-modul bergradasi, yakni memenuhi $\varphi(A_i) \subset B_i$. \emph{Ideal bergradasi} berarti ideal dalam $A$ berbentuk $\mathfrak{a} = \bigoplus _i (\mathfrak{a} \cap A_i)$; jika $\mathfrak{a} \neq A$, maka $A/\mathfrak{a} = \bigoplus_i A_i/\mathfrak{a} \cap A_i$ tetap merupakan aljabar bergradasi.
	\end{itemize}
	Jika $(I,+) \subset (\Z,+)$, kita cukup menyebut objek-objek tersebut bergradasi.
\end{definition}

\begin{lemma}\label{prop:graded-ideal-submodule}
	Dalam $I$-aljabar bergradasi (atau modul bergradasi), suatu ideal dua sisi $\mathfrak{a}$ (atau submodul) bergradasi jika dan hanya jika dapat dibangkitkan oleh unsur-unsur homogen.
\end{lemma}
\begin{proof}
	Kita hanya membahas kasus ideal. Menurut definisi, ideal bergradasi tentu dibangkitkan oleh unsur-unsur homogen. Untuk membuktikan arah sebaliknya, cukup tunjukkan bahwa jika $\mathfrak{a}$ mempunyai pembangkit homogen $\{a_s : s \in S \}$, maka $\mathfrak{a} = \sum_i \mathfrak{a} \cap A_i$. Menurut definisi pembangkit, unsur-unsur $\mathfrak{a}$ merupakan kombinasi $R$-linear dari unsur berbentuk $u a_s v$, dengan $u,v \in A$. Dengan menguraikan $u, v$ lebih lanjut ke dalam suku-suku homogen, kita dapat mengasumsikan semua $u a_s v$ tersebut homogen. Terbukti.
\end{proof}

Setiap $R$-aljabar $A$ dapat diberi struktur $I$-bergradasi trivial: tetapkan $A_0 = A$ dan $A_i = \{0\}$ untuk indeks lainnya. Untuk $A = R$, inilah satu-satunya pilihan. Karena itu, definisi modul bergradasi dapat diperluas secara wajar sebagai berikut. Jika $R$-modul bergradasi $M$ sekaligus merupakan modul kiri atas $R$-aljabar $A$, maka $M$ disebut modul bergradasi atas $A$ apabila $A_i \cdot M_j \subset M_{i+j}$ untuk semua $i,j \in I$.

Pendekatan yang lebih rapi ialah memandang aljabar bergradasi $A$ sebagai struktur yang ditumpangkan pada modul bergradasi. Ini sama dengan mensyaratkan dalam definisi aljabar melalui panah, Definisi \ref{def:algebra-mod-diagram}, bahwa $A$ merupakan $I$-modul bergradasi dan bahwa perkalian $A \otimes A \to A$ serta satuan $R \to A$ keduanya merupakan homomorfisme $I$-modul bergradasi. Mari kita lihat beberapa contoh awal.

\begin{itemize}
	\item Gelanggang polinomial $R[X_0, \ldots, X_n]$ adalah $\Z_{\geq 0}$-aljabar bergradasi, dengan
		\[ R[X_0, \ldots, X_n]_i = \bigoplus_{|\bm{a}|=i} R X_0^{a_0}  \cdots X_n^{a_n}, \]
		sehingga unsur homogen berderajat $i$ tepat merupakan polinomial homogen berderajat $i$ atas $R$. Ideal bergradasi dalam kasus ini juga disebut ideal homogen; dengan beberapa pengecualian kecil, dalam geometri aljabar ideal-ideal itu bersesuaian dengan subskema tertutup dalam ruang projektif $\mathbb{P}^n$.
	\item Misalkan $\Gamma$ monoid komutatif. Gelanggang monoid dari Definisi \ref{def:monoidal-ring}, yaitu $R[\Gamma] = \bigoplus_{\gamma \in \Gamma} R\gamma$, membentuk $\Gamma$-aljabar bergradasi; ini hampir merupakan pengulangan definisi.
	\item Semua bentuk diferensial bernilai $\CC$ pada manifold diferensiabel mulus $X$, dengan operasi $(\omega, \eta) \mapsto \omega \wedge \eta$, membentuk $\CC$-aljabar yang dinotasikan $A(X) = \bigoplus_{k=0}^{\dim X} A^k(X)$. Di sini $A^k(X)$ adalah ruang vektor atas $\CC$ yang terdiri atas bentuk diferensial berderajat $k$; teori dasarnya dapat dilihat dalam \cite[\S 3.2]{ChCh}. Ingat bahwa $A^0(X) = C^\infty(X)$, sedangkan setiap unsur $\omega$ dari $A^k(X)$ dapat ditulis secara unik dalam koordinat lokal $x_1, \ldots, x_n$ sebagai
	\[ \omega = \sum_{1 \leq i_1 < \cdots < i_k \leq n} f_{i_1, \ldots, i_k} \dd x_{i_1} \wedge \cdots \wedge \dd x_{i_k}, \quad f_{i_1, \ldots, i_k} \in C^\infty(X) \]
	dengan demikian $A(X)$ merupakan aljabar bergradasi. Perkaliannya tidak komutatif, tetapi ``tidak jauh dari komutatif'', karena dalam geometri berlaku rumus yang dikenal
	\[  \omega \wedge \eta = (-1)^{ab} \eta \wedge \omega, \quad \omega \in A^a(X), \; \eta \in A^b(X). \]
	Sifat penting lain aljabar $A(X)$ ialah adanya operasi diferensial eksterior $d: A(X) \to A(X)$.
\end{itemize}
Aljabar $A(X)$ dalam contoh terakhir mengandung struktur aljabar yang kaya dan merupakan perangkat yang tak tergantikan dalam geometri, terutama berkat karya Élie Cartan. Dari rumus pertukaran perkalian pada $A(X)$, kita abstraksikan definisi berikut.

\begin{definition}\label{def:anti-commutation}\index{daishu!$\epsilon$-komutatif}
	Berikan monoid komutatif $(I, +)$ dan homomorfisme aditif $\epsilon: I \to \Z/2\Z$. Jika aljabar bergradasi $A$ memenuhi
	\[ xy = (-1)^{\epsilon(\deg x) \epsilon(\deg y)} yx, \quad x, y \in A: \;\text{unsur homogen} \]
	maka perkalian $A$ disebut $\epsilon$-komutatif.
\end{definition}
Untuk $\epsilon=0$, kita kembali memperoleh aljabar bergradasi komutatif. Pilihan yang paling lazim tentu $I \subset \Z$ dan $\epsilon: \Z \to \Z/2\Z$ adalah homomorfisme hasil bagi; dalam kasus ini, $\epsilon$-komutatif sering disebut \emph{antikomutatif}. Kadang-kadang kita juga memberlakukan sifat
\begin{gather}\label{eqn:strong-anti-commutation}
	\left[ \epsilon(\deg x) = 1 \in \Z/2\Z \right] \implies x^2 = 0.
\end{gather}
Jika $2 \in R^\times$, sifat $\epsilon$-komutatif menyiratkan $\epsilon(\deg x)=1 \implies x^2 = -x^2 \implies x^2=0$, sehingga syarat \eqref{eqn:strong-anti-commutation} berlaku otomatis. Dalam penerapan seperti geometri diferensial, biasanya $R = \CC, \R$, sehingga tidak ada perbedaan.

Dari definisi mudah diperiksa bahwa hasil kali tensor $I$-aljabar bergradasi $A$ dan $B$, dengan $(A \otimes B)_i = \bigoplus_{j+k=i} A_j \otimes B_k$, tetap merupakan $I$-aljabar bergradasi, dan kendala komutativitas $c(A,B): A \otimes B \rightiso B \otimes A$ adalah isomorfisme aljabar bergradasi. Tentu saja, hasil serupa berlaku untuk hasil kali tensor tak hingga dari aljabar bergradasi. Dalam kasus aljabar bergradasi antikomutatif, misalnya $A(X)$, perkalian pada $A \otimes B$ perlu sedikit ``dipuntir'' untuk memperoleh sifat yang diharapkan. Kita hendak memahami mekanisme ini dalam kerangka yang seragam.

Sudut pandang yang paling ampuh untuk menangani persoalan tersebut ialah memperkenalkan sebuah struktur kepang pada kategori $I$-modul bergradasi $R\dcate{Mod}_I$ (Definisi \ref{def:braiding}), yaitu isomorfisme alami $c_\epsilon(M, N): M \otimes N \rightiso N \otimes M$ berikut:
\begin{align*}
	c_\epsilon(M,N): M_i \otimes N_j & \longrightiso N_j \otimes M_i \\
	x \otimes y & \longmapsto (-1)^{\epsilon(\deg x)\epsilon(\deg y)} y \otimes x, \quad i,j \in I.
\end{align*}
Struktur ini juga disebut struktur kepang Koszul\index{bianjiegou!Koszul}. Mudah diperiksa bahwa ia memenuhi semua syarat dalam Definisi \ref{def:braiding}, bahkan juga memenuhi sifat simetris
\begin{gather}\label{eqn:Koszul-braiding-symm}
	c_\epsilon(A,B) c_\epsilon(B, A) = \identity.
\end{gather}
Selain itu, Definisi \ref{def:anti-commutation} dapat ditulis ulang dalam bahasa kategori monoidal sebagai
\[ A:\; \epsilon\text{-komutatif} \iff \mu_A \circ c_\epsilon(A,A) = \mu_A \]
dengan $\mu_A: A \otimes A \to A$ menyatakan perkalian pada $A$.

\begin{definition-theorem}[Hukum tanda Koszul]\label{def:Koszul-sign-alg}
	Tetapkan $I$ dan $\epsilon: I \to \Z/2\Z$, serta misalkan $A$ dan $B$ merupakan $I$-aljabar bergradasi. Definisikan perkalian pada hasil kali tensor $A \otimes B$ sehingga untuk unsur-unsur homogen berlaku
	\[ (a \otimes b)(a' \otimes b') = (-1)^{\epsilon(\deg b)\epsilon(\deg a')} aa' \otimes bb' , \]
	maka
	\begin{compactenum}[(i)]
		\item $A \otimes B$ dengan perkalian tersebut menjadi $I$-aljabar bergradasi;
		\item homomorfisme alami $\iota_A: A \to A \otimes B$ dan $\iota_B: B \to A \otimes B$ merupakan homomorfisme aljabar bergradasi;
		\item jika $A$ dan $B$ keduanya $\epsilon$-komutatif, maka demikian pula $A \otimes B$;
		\item struktur kepang Koszul $c_\epsilon(A, B): A \otimes B \rightiso B \otimes A$ merupakan isomorfisme aljabar bergradasi.
	\end{compactenum}
\end{definition-theorem}
Dengan kata lain, pada setiap komponen bergradasi $(A_i \otimes B_j) \otimes (A_k \otimes B_l) \to A_{i+k} \otimes B_{j+l}$ dari perkalian hasil kali tensor semula, kita memuntirnya dengan tanda $(-1)^{\epsilon(j)\epsilon(k)}$. Pembuktiannya hanya merupakan pemeriksaan langsung dari definisi; pokoknya ialah memahami gagasan di balik definisi. Andaikan $A$ dan $B$ keduanya $\epsilon$-komutatif. Syarat apa sekurang-kurangnya diperlukan agar $A \otimes B$ menjadi $\epsilon$-komutatif? Pernyataan (ii) merupakan tuntutan yang alami dan menentukan aturan perkalian $(a \otimes 1)(a' \otimes 1) = aa' \otimes 1$ serta $(1 \otimes b)(1 \otimes b') = 1 \otimes bb'$. Dalam kasus umum, hukum asosiatif menyiratkan
\[ (a \otimes b)(a' \otimes b') = (a \otimes 1) \left((1 \otimes b)(a' \otimes 1)\right) (1 \otimes b'). \]
Jika $A \otimes B$ bersifat $\epsilon$-komutatif, harus berlaku $(1 \otimes b)(a' \otimes 1) = (-1)^{\epsilon(\deg b)\epsilon(\deg a)} (a' \otimes 1)(1 \otimes b)$. Substitusi ke langkah sebelumnya menghasilkan tepat aturan perkalian yang dipuntir oleh $\epsilon$ dalam definisi. Jadi, di bawah batasan (ii), hukum tanda Koszul bukan hanya alami, tetapi juga niscaya bagi $I$-aljabar bergradasi yang $\epsilon$-komutatif. Hal ini dapat diringkas dengan rapi dalam bahasa kategori.

\begin{proposition}\label{prop:graded-alg-cocomplete}
	Misalkan $R\dcate{CAlg}^\epsilon_I$ kategori semua $I$-aljabar bergradasi yang $\epsilon$-komutatif. Maka hasil kali tensor $\otimes$, bersama hukum tanda Koszul, memberikan koproduk dalam $R\dcate{CAlg}^\epsilon_I$. Selain itu, $R\dcate{CAlg}^\epsilon_I$ kokomplet. Jika $A$ dan $B$ keduanya memenuhi \eqref{eqn:strong-anti-commutation}, maka $A \otimes B$ juga memenuhinya.
\end{proposition}
Kategori $I$-aljabar bergradasi komutatif $R\dcate{CAlg}^0_I$ biasanya disingkat menjadi $R\dcate{CAlg}_I$.
\begin{proof}
	Bagian pertama merupakan perluasan alami dari Korolari \ref{prop:tensor-alg-cocomplete}; konstruksi koproduk hingga pada dasarnya telah dijelaskan di atas, dan limit terarah ke atas dapat digunakan lagi untuk menangani hasil kali tensor tak hingga. Untuk membuktikan bahwa $R\dcate{CAlg}^\epsilon_I$ kokomplet, bagi setiap pasangan morfisme $f,g: A \to B$ cukup dibangun koekualiser $B \to \Coker(f,g)$. Perhatikan bahwa $\{a \in A:f(a)-g(a) \}$ membangkitkan ideal bergradasi $I$, sehingga $B/I$ adalah objek yang dicari.
	
	Sekarang andaikan $A$ dan $B$ memenuhi \eqref{eqn:strong-anti-commutation}. Tinjau dalam $A \otimes B$ unsur homogen $x = \sum_{i=1}^n x_i \otimes y_i$. Tuliskan $p_i := \epsilon(\deg x_i)$ dan $q_i := \epsilon(\deg y_i)$, serta andaikan $\forall i,\; p_i+q_i=1$. Hitung
	\begin{align*}
		x^2 & = \left( \sum_{i=1}^n x_i \otimes y_i \right)^2 = \sum_{i,j} (-1)^{q_i p_j} x_i x_j \otimes y_i y_j \\
		& = \left( \sum_{i=j} + \sum_{i \neq j} \right) \cdots
		= 0 + \sum_{i \neq j} (-1)^{q_i p_j} x_i x_j \otimes y_i y_j.
	\end{align*}
	Untuk setiap $1 \leq i,j \leq n$,
	\begin{align*}
		(-1)^{q_i p_j} x_i x_j \otimes y_i y_j & = (-1)^{q_i p_j + p_i p_j + q_i q_j} x_j x_i \otimes y_j y_i \\
		& = (-1)^{p_i q_j + q_i p_j + p_i p_j + q_i q_j} \cdot (-1)^{q_j p_i} x_j x_i \otimes y_j y_i.
	\end{align*}
	Karena dalam $\Z/2\Z$ berlaku $q_j p_i + q_i p_j + p_i p_j + q_i q_j = (p_i + q_i)(p_j + q_j) = 1$, suku-suku dalam jumlah di atas berubah tanda dan saling meniadakan ketika $i \neq j$ dipertukarkan. Maka, sebagaimana diinginkan, $x^2=0$.
\end{proof}

Kembali ke asal geometrisnya, mudah dilihat bahwa hukum tanda Koszul bersama pemetaan $\omega \otimes \eta \mapsto \omega \wedge \eta$ hampir mengidentifikasi $A(X) \otimes A(Y)$ dengan $A(X \times Y)$, sehingga konstruksi ini wajar. Cukup pilih koordinat lokal $x_1, \ldots, x_n$ pada $X$ dan $y_1, \ldots, y_m$ pada $Y$; gabungannya menjadi koordinat lokal pada $X \times Y$. Perkalian bentuk diferensial mengharuskan bahwa dalam $A(X \times Y)$ berlaku
\[ \left( \dd y_{i_1} \wedge \cdots \wedge \dd y_{i_k} \right) \wedge \left( \dd  x_{j_1} \wedge \cdots \wedge \dd x_{j_l} \right) = (-1)^{kl} \dd x_{j_1} \wedge \cdots \wedge \dd x_{j_l} \wedge \dd y_{i_1} \wedge \cdots \wedge \dd y_{i_k} \]
Ini tepat merupakan hukum tanda Koszul. Kata ``hampir'' diperlukan karena $C^\infty(X) \otimes C^\infty(Y) \neq C^\infty(X \times Y)$; untuk memperoleh persamaan, harus dipakai hasil kali tensor lengkap $\hat{\otimes}$ dari ruang vektor topologis.

Dari sudut pandang kategori monoidal simetris, pandang aljabar bergradasi sebagai struktur yang ditumpangkan pada modul bergradasi. Hukum tanda Koszul berarti, pada $A \otimes B$, mendefinisikan perkalian $\mu_{A \otimes B}$ sebagai morfisme komposit berikut (dengan kendala asosiativitas dihilangkan):
\[ A \otimes B \otimes A \otimes B \xrightarrow{\identity_A \otimes c_\epsilon(B, A) \otimes \identity_B} A \otimes A \otimes B \otimes B \xrightarrow{\mu_A \otimes \mu_B} A \otimes B; \]
Untuk $\epsilon=0$ dan $c_0(B,A) = c(B, A)$, diperoleh versi tanpa puntiran. Sifat $\epsilon$-komutatif dari aljabar bergradasi $A$ juga dapat dicirikan oleh $A = A^{\epsilon-\text{op}}$, dengan aljabar lawan-$\epsilon$ $A^{\epsilon-\text{op}}$ dibentuk dengan menggunakan $\mu_A \circ c_\epsilon(A,A)$ sebagai pengganti $\mu_A$. Pembaca dapat memeriksa bahwa $A^{\epsilon-\text{op}}$ memang suatu aljabar.

Ringkasnya, dengan meninjau struktur kepang simetris yang berbeda-beda pada kategori monoidal yang sama $(R\dcate{Mod}_I, \otimes)$, kita memperoleh berbagai versi komutativitas bagi aljabar bergradasi dan berbagai perkalian pada hasil kali tensor. Gagasan ini bukanlah contoh yang terisolasi dalam matematika modern.

Definisi--Teorema \ref{def:Koszul-sign-alg} pada akhirnya hanya memuat sifat-sifat aljabar dasar yang mudah diperiksa langsung. Namun, setelah mencapai sudut pandang umum ini, pemeriksaan dalam kerangka kategori monoidal simetris mungkin lebih menarik dan juga memudahkan perumuman selanjutnya. Sebagai contoh, untuk sifat (iv) kita perlu membuktikan bahwa diagram
\[\begin{tikzcd}[column sep=large, row sep=large]
	A \otimes B \otimes A \otimes B \arrow[r, "{c_\epsilon(A,B) \otimes c_\epsilon(A,B)}" inner sep=1em] \arrow[d, "{\identity_A \otimes c_\epsilon(B,A) \otimes \identity_B}"'] & B \otimes A \otimes B \otimes A \arrow[d, "{\identity_B \otimes c_\epsilon(A,B) \otimes A}"] \\
	A \otimes A \otimes B \otimes B \arrow[r, "{c_\epsilon(A \otimes A, B \otimes B)}" inner sep=1em] \arrow[d, "{\mu_A \otimes \mu_B}"'] &  B \otimes B \otimes A \otimes A \arrow[d, "{\mu_B \otimes \mu_A}"] \\
	A \otimes B \arrow[r, "{c_\epsilon(A,B)}"'] & B \otimes A
\end{tikzcd}\]
komutatif, dengan $\mu_A$ dan $\mu_B$ menyatakan perkalian aljabar dan kendala asosiativitas dihilangkan dari diagram. Komutativitas bagian bawah mengikuti dari fungtorialitas $c_\epsilon(-,-)$. Bagian atas langsung terlihat dari sudut pandang kepang (ingat \S\ref{sec:braiding}):
\begin{equation*}\begin{tikzpicture}[baseline=(braid)]
	\braid[style strands={1}{black}, style strands={2}{black}, style strands={3}{black}, style strands={4}{black},
	floor command={
		\draw[dashed] (\floorsx,\floorsy) -- (\floorex,\floorsy);
	}] (braid)  s_2 | s_1-s_3;
	\node[above] at (braid-1-s) {$B$};
	\node[below] at (braid-1-e) {$B$};
	\node[above] at (braid-2-s) {$B$};
	\node[below] at (braid-2-e) {$B$};
	\node[above] at (braid-3-s) {$A$};
	\node[below] at (braid-3-e) {$A$};
	\node[above] at (braid-4-s) {$A$};
	\node[below] at (braid-4-e) {$A$};
\end{tikzpicture} \; = \;
\begin{tikzpicture}[baseline=(braid), yscale=0.7]
	\braid[style strands={1}{black}, style strands={2}{black}, style strands={3}{black}, style strands={4}{black} ] (braid) s_2 s_1-s_3 s_2 s_2^{-1};
	\node[above] at (braid-1-s) {$B$};
	\node[below] at (braid-1-e) {$B$};
	\node[above] at (braid-2-s) {$B$};
	\node[below] at (braid-2-e) {$B$};
	\node[above] at (braid-3-s) {$A$};
	\node[below] at (braid-3-e) {$A$};
	\node[above] at (braid-4-s) {$A$};
	\node[below] at (braid-4-e) {$A$};
\end{tikzpicture} \; \xlongequal{\because\,\eqref{eqn:Koszul-braiding-symm}} \;
\begin{tikzpicture}[baseline=(braid), yscale=0.7]
	\braid[style strands={1}{black}, style strands={2}{black}, style strands={3}{black}, style strands={4}{black},
	floor command={
		\draw[dashed] (\floorsx,\floorsy) -- (\floorex,\floorsy);
	}] (braid) s_2 s_1-s_3 s_2 | s_2;
	\node[above] at (braid-1-s) {$B$};
	\node[below] at (braid-1-e) {$B$};
	\node[above] at (braid-2-s) {$B$};
	\node[below] at (braid-2-e) {$B$};
	\node[above] at (braid-3-s) {$A$};
	\node[below] at (braid-3-e) {$A$};
	\node[above] at (braid-4-s) {$A$};
	\node[below] at (braid-4-e) {$A$};
\end{tikzpicture}\end{equation*}
Pembuktian formal mengharuskan kita terlebih dahulu menguraikan $c_\epsilon(A \otimes A, B \otimes B)$ memakai \eqref{eqn:hexagon-axiom-1} dan \eqref{eqn:hexagon-axiom-2} menjadi bentuk bagian atas kepang paling kanan, tetapi hal ini juga tidak sulit.

\begin{remark}
	Fungtor $P_{R \to S}$ dan $\mathcal{F}_{R \to S}$ mempunyai perumuman yang jelas ke kasus $I$-aljabar bergradasi dan modul bergradasi, beserta relasi adjoin yang terkait. Sebagai contoh, untuk perubahan basis $P_{R \to S}$, jika $S$ diberi gradasi trivial $S_0 = S$, maka $P_{R \to S} = - \otimes S$ secara alami meluas menjadi fungtor $R\dcate{Alg}_I \to S\dcate{Alg}_I$. Perluasan ke aljabar $\epsilon$-simetris juga langsung.
\end{remark}

\section{Aljabar Tensor}
Tetapkan sebuah gelanggang komutatif $R$. Dalam \S\ref{sec:module-tensor-prod} kita telah membahas hasil kali tensor modul secara sistematis; dalam situasi komutatif yang sedang kita hadapi, persoalannya lebih sederhana. Selanjutnya tuliskan $\otimes = \otimes_R$. Kita dapat meninjau hasil kali tensor beberapa modul $M_1 \otimes (M_2 \otimes (\cdots \otimes M_n )\cdots)$; berkat kendala asosiativitas, penempatan tanda kurung tidak menjadi masalah. Cara untuk membereskannya sekaligus ialah meniru bukti Proposisi \ref{prop:tensor-assoc}: untuk setiap $R$-modul $A$, tinjau $R$-modul
\[ \text{Mul}(M_1, \cdots, M_n; A) := \left\{ \text{pemetaan multilinear}\; B: M_1 \times \cdots \times M_n \to A \right\} \]
Di sini $B$ disebut pemetaan multilinear apabila ia $R$-linear dalam setiap peubah $x_i \in M_i$.\index{duochongxianxingyingshe@Pemetaan multilinear (multilinear map)}

Data yang hendak kita definisikan adalah $R$-modul $M_1 \otimes \cdots \otimes M_n$ beserta pemetaan multilinear $M_1 \times \cdots \times M_n \to M_1 \otimes \cdots \otimes M_n$, yang terakhir ditulis $(x_1, \cdots, x_n) \mapsto x_1 \otimes \cdots \otimes x_n$. Keduanya dicirikan oleh sifat universal
\begin{equation}\label{eqn:Mul-tensor}\begin{aligned}
	\Hom(M_1 \otimes \cdots \otimes M_n, \bullet) & \stackrel{\sim}{\longrightarrow} \text{Mul}(M_1, \cdots, M_n; \bullet) \\
	\varphi & \longmapsto [ (x_1, \ldots, x_n) \mapsto \varphi(x_1 \otimes \cdots \otimes x_n) ]
\end{aligned}\end{equation}
Di sini tetap terdapat kendala asosiativitas multipeubah
\[ (M_1 \otimes \cdots \otimes M_n) \otimes (M_{n+1} \otimes \cdots \otimes M_m) \rightiso M_1 \otimes \cdots \otimes M_m \]
yang memenuhi sifat fungtorial yang semestinya. Hasil kali tensor ke-$n$ dari modul $M$ didefinisikan sebagai
\begin{align*}
	T^n(M) & := \underbracket{M \otimes \cdots \otimes M}_{n \text{ salinan}}, \quad n \geq 1, \\
	T^0(M) & := R.
\end{align*}
Kendala asosiativitas menginduksi homomorfisme alami $\mu_{i,j}: T^i M \otimes T^j M \rightiso T^{i+j} M$.

\begin{definition}[Aljabar tensor]\index[sym1]{$T(M)$, $T^n(M)$}\index{zhangliangdaishu@Aljabar tensor (tensor algebra)}
	\emph{Aljabar tensor} dari $R$-modul $M$ didefinisikan sebagai $T(M) := \bigoplus_{n=0}^\infty T^n(M)$, dengan perkalian dan unsur satuan masing-masing diberikan oleh semua $\mu_{i,j}$ dan $R = T^0(M) \hookrightarrow T(M)$. Ia dilengkapi monomorfisme alami $R$-modul $M = T^1(M) \hookrightarrow T(M)$. Selain itu, $T(M)$ juga merupakan aljabar bergradasi dalam arti Definisi \ref{def:graded-module-alg}, dengan bagian homogen berderajat $k$ tepat $T^k(M)$.
\end{definition}
Jika definisi perkaliannya dijabarkan pada tingkat unsur, ia tidak lain adalah hasil kali tensor
\[ (x_1 \otimes \cdots \otimes x_n) \cdot (y_1 \otimes \cdots \otimes y_m) = x_1 \otimes \cdots \otimes x_n \otimes y_1 \otimes \cdots \otimes y_m. \]
Sifat-sifat yang diperlukan, seperti hukum asosiatif, pada akhirnya berasal dari fungtorialitas kendala asosiativitas. Perhatikan bahwa $M = T^1(M)$ membangun $T(M)$. Karena setiap $T^k(\cdot)$ merupakan fungtor, mudah dibuktikan bahwa $T(\cdot)$ juga demikian: jika $M \to N$ adalah homomorfisme modul, maka terdapat tepat satu homomorfisme aljabar $T(M) \to T(N)$ yang membuat diagram
$\begin{tikzcd}[row sep=small, column sep=small]
	M \arrow[r] \arrow[d] & N \arrow[d] \\
	T(M) \arrow[r] & T(N)
\end{tikzcd}$
komutatif.

Sifat universal berikut menunjukkan bahwa $(T(M), M \to T(M))$ dapat disebut ``aljabar bebas'' di atas modul $M$.
\begin{theorem}\label{prop:tensor-alg-universal}
	Aljabar tensor memenuhi sifat universal berikut: untuk setiap $R$-aljabar $A$ beserta homomorfisme $R$-modul $f: M \to A$, terdapat tepat satu homomorfisme $R$-aljabar $\varphi: T(M) \to A$ yang membuat diagram berikut komutatif:
	\[\begin{tikzcd}[row sep=tiny]
		& T(M) \arrow[dd, "\exists! \;\varphi"] \\
		M \arrow[ru] \arrow[rd, "f"'] & \\
		& A
	\end{tikzcd}\]
	Dengan kata lain, terdapat relasi adjoin antarfungtor $\Hom_{R\dcate{Alg}}(T(-), -) \simeq \Hom_{R\dcate{Mod}}(-, U(-))$, dengan $U: R\dcate{Alg} \to R\dcate{Mod}$ fungtor pelupa.
\end{theorem}
\begin{proof}
	Homomorfisme $\varphi: T(M) \to A$ ditentukan oleh pembatasannya pada setiap suku jumlah langsung $T^n(M)$; menurut sifat universal, untuk $n \geq 1$ pembatasan itu selanjutnya ditentukan oleh pemetaan multilinear
	\begin{align*}
		M \times \cdots \times M & \longrightarrow A \\
		(x_1, \cdots, x_n) & \longmapsto \varphi(x_1 \otimes \cdots \otimes x_n)
	\end{align*}
	Homomorfisme aljabar dalam diagram harus memenuhi $\varphi(x_1 \otimes \cdots \otimes x_n) = f(x_1) \cdots f(x_n)$; ruas kanan adalah pemetaan multilinear dalam $(x_1, \ldots, x_n)$. Sifat universal karenanya menentukan $\varphi$ secara tunggal. Untuk membuktikan eksistensi, cukup balik penalaran di atas untuk membangun homomorfisme modul $\varphi: T(M) \to A$, lalu periksa sifat multiplikatifnya secara langsung.
\end{proof}

Karena fungtor adjoin telah muncul, mari kita ingat kembali prosedur baku teori kategori. Isomorfisme $\Hom_{R\dcate{Alg}}(T(-), -) \simeq \Hom_{R\dcate{Mod}}(-, U(-))$ memetakan $\identity_{T(M)}: T(M) \to T(M)$ di ruas kiri ke homomorfisme $R$-modul $M \to UT(M) = T(M)$; inilah morfisme ``unit'' yang menentukan isomorfisme adjoin secara tunggal (Definisi \ref{def:adjunction-unit-counit}). Jadi penafsiran melalui adjoin sekaligus mencakup konstruksi $T(M)$ dan $M \to T(M)$.

\begin{lemma}\label{prop:tensor-alg-basechange}
	Konstruksi aljabar tensor komutatif dengan perubahan basis: untuk homomorfisme gelanggang $R \to S$, terdapat tepat satu isomorfisme aljabar $S$ bergradasi $\psi_M: T(M \otimes S) \rightiso T(M) \otimes S$ yang membuat diagram
	\[\begin{tikzcd}
		& M \otimes S \arrow[ld] \arrow[rd, "{[M \to T(M)] \otimes S}"] & \\
		T(M \otimes S) \arrow[rr, "\sim", "\psi_M"'] & & T(M) \otimes S
	\end{tikzcd}\]
	komutatif di dalam $S\dcate{Mod}$. Dengan demikian terdapat isomorfisme fungtor $T(- \otimes S) \rightiso T(-) \otimes S$.
\end{lemma}
\begin{proof}
	Karena $M \rightiso T^1(M)$ membangun $R$-aljabar $T(M)$, ketunggalan dan sifat bergradasi dari isomorfisme $\psi_M$ sudah jelas. Ada berbagai cara untuk membuktikan eksistensinya; pemeriksaan langsung pada tingkat himpunan unsur dan pemetaan mungkin yang paling cepat, tetapi kita ingin menonjolkan peran sentral fungtor adjoin. Maka tinjau fungtor perubahan basis $P_{R \to S}: R\dcate{Mod} \xrightarrow{- \otimes S} S\dcate{Mod}$ dan fungtor tarik-balik $\mathcal{F}_{R \to S}: S\dcate{Mod} \to R\dcate{Mod}$; versi keduanya untuk $R$- dan $S$-aljabar masing-masing ditulis $P_{R \to S}^a$, $\mathcal{F}^a_{R \to S}$. Korolari \ref{prop:IP-vs-F} dan Proposisi \ref{prop:IP-vs-F-alg} memberikan pasangan adjoin $(P_{R \to S}, \mathcal{F}_{R \to S})$ dan $(P_{R \to S}^a, \mathcal{F}_{R \to S}^a)$; selain itu, Teorema \ref{prop:tensor-alg-universal} memberikan pasangan adjoin $(T, U)$. Demi menyederhanakan notasi, isomorfisme yang membentuk pasangan-pasangan adjoin tersebut tidak dituliskan.

	Dengan mengomposisikan pasangan-pasangan adjoin menurut Proposisi \ref{prop:adjunction-composition}, kita memperoleh bahwa $(T P_{R \to S}, \mathcal{F}_{R \to S} U)$ dan $(P_{R \to S}^a T, U \mathcal{F}_{R \to S}^a)$ juga merupakan pasangan adjoin. Jelas bahwa $\mathcal{F}_{R \to S} U = U \mathcal{F}_{R \to S}^a$, sehingga ketunggalan dalam Proposisi \ref{prop:adjunction-uniqueness} memberikan $\psi: TP_{R \to S} \rightiso P_{R \to S}^a T$. Dengan meninjau unit (Definisi \ref{def:adjunction-unit-counit}), diperoleh diagram komutatif alami
	\[\begin{tikzcd}[row sep=small]
		{} & M \arrow[ld] \arrow[rd] & \\
		\mathcal{F}_{R \to S} U T P_{R \to S} (M) \arrow[rr, "\sim", "{\mathcal{F}_{R \to S}U \psi_M}"'] & & \mathcal{F}_{R \to S} U P_{R \to S}^a T(M)
	\end{tikzcd}\]
	dan diagram ini juga menentukan $\psi_M$ secara tunggal (sebab $U$ dan $\mathcal{F}_{R \to S}$ keduanya setia). Dengan menggunakan pasangan adjoin $(P_{R \to S}, \mathcal{F}_{R \to S})$ beserta kealamiannya, pernyataan ini ekuivalen dengan komutativitas diagram
	\[\begin{tikzcd}[row sep=small]
		{} & P_{R \to S}(M) \arrow[ld] \arrow[rd] & \\
		U T P_{R \to S} (M) \arrow[rr, "\sim", "{U \psi_M}"'] & & U P_{R \to S}^a T(M)
	\end{tikzcd}\]
	Dengan menjabarkan definisi fungtor-fungtor dalam diagram, kita memperoleh diagram komutatif yang dinyatakan dalam lemma. Dari sini juga tampak bahwa $\psi_M$ mempertahankan gradasi.
\end{proof}

\begin{corollary}\label{prop:tensor-alg-limit}
	Konstruksi aljabar tensor $T(-): R\dcate{Mod} \to R\dcate{Alg}_{\Z}$ mempertahankan setiap $\varinjlim$: terdapat isomorfisme alami aljabar bergradasi $\varinjlim_i T(M_i) \rightiso T(\varinjlim_i M_i)$. Epimorfisme modul $M \twoheadrightarrow M/N$ dipetakan ke $T(M) \to T(M)/\lrangle{N} \simeq T(M/N)$, dengan $\lrangle{N}$ ideal bergradasi yang dibangun oleh $N$.
\end{corollary}
\begin{proof}
	Karena $T$ memiliki fungtor adjoin kanan $U$, dipertahankannya $\varinjlim$ merupakan konsekuensi formal Teorema \ref{prop:adjuncion-limit}. Epimorfisme $M \to M/N$ dapat dipandang di dalam $R\dcate{Mod}$ sebagai koekualiser $M \to \Coker(\iota,0)$, dengan $\iota: N \hookrightarrow M$. Ia dipetakan oleh $T$ ke $T(M) \to \Coker(T(\iota), T(0))$. Untuk memperoleh isomorfisme alami $\Coker(T(\iota), T(0)) \simeq T(M)/\lrangle{N}$, cukup ingat konstruksi koekualiser di dalam $R\dcate{Alg}_{\Z}$ (Proposisi \ref{prop:graded-alg-cocomplete}) dan perhatikan bahwa $T(0): T(N) \to T(M)$ menginduksi $R = T^0(N) \hookrightarrow T^0(M)$ serta nol pada $T^{>0}(N)$.
\end{proof}

\section{Aljabar simetris dan aljabar eksterior}
Gunakan notasi dari bagian sebelumnya.
\begin{definition}[Aljabar simetris dan aljabar eksterior]\label{def:Sym-wedge} \index{duichengdaishu@Aljabar simetris (symmetric algebra)} \index{waichengdaishu@Aljabar eksterior (exterior algebra)} \index[sym1]{Sym(M)@$\Sym(M)$, $\Sym^n(M)$} \index[sym1]{Lambda(M)@$\bigwedge(M)$, $\bigwedge^n(M)$}
	Definisikan ideal dua sisi bergradasi dari $T(M)$ melalui pembangun-pembangun homogen
	\begin{align*}
		I_{\Sym}(M) & := \lrangle{ x \otimes y - y \otimes x : x,y \in M}, \\
		I_{\bigwedge}(M) & := \lrangle{x \otimes x : x \in M}.
	\end{align*}
	Aljabar-aljabar hasil bagi yang bersesuaian
	\begin{align*}
		\Sym(M) & := T(M)/I_{\Sym}(M), \\
		\bigwedge(M) & := T(M)/I_{\bigwedge}(M)
	\end{align*}
	masing-masing disebut \emph{aljabar simetris} dan \emph{aljabar eksterior} dari $M$.
\end{definition}

Ketika $R = \CC$ dan $M$ adalah ruang Hilbert, $\Sym(M)$ merupakan ruang Fock (bosonik) yang lazim dalam fisika kuantum, sedangkan $\bigwedge(M)$ menghasilkan ruang Fock (fermionik); peralihan dari $M$ ke ruang Fock-nya merupakan salah satu langkah dalam apa yang disebut kuantisasi kedua.

Kita mulai dengan beberapa pengamatan dasar.
\begin{itemize}
	\item Lemma \ref{prop:graded-ideal-submodule} memastikan bahwa $I_{\Sym}(M) = \bigoplus_n I_{\Sym}^n(M)$ dan $I_{\bigwedge}(M) = \bigoplus_n I_{\bigwedge}^n(M)$ keduanya ideal bergradasi. Karena itu, $\Sym(M) = \bigoplus_{n \geq 0} \Sym^n(M)$ dan $\bigwedge(M) = \bigoplus_{n \geq 0} \bigwedge^n(M)$ secara alami menjadi aljabar bergradasi. Karena ideal-ideal ini dibangun oleh unsur homogen berderajat dua, dari $M \to T(M)$ kita tetap memperoleh monomorfisme $M \to \Sym(M)$ dan $M \to \bigwedge(M)$.
	\item Setiap homomorfisme modul $M \to N$ menginduksi $I_{\Sym}(M) \to I_{\Sym}(N)$ dan $I_{\bigwedge}(M) \to I_{\bigwedge}(N)$; dengan demikian, $\Sym(\cdot)$ dan $\bigwedge(\cdot)$ keduanya merupakan fungtor.
	\item Perkalian dalam aljabar simetris lazim ditulis $(x,y) \mapsto xy$; jika $x, y \in M$, maka $xy=yx$. Karena $M$ membangun $\Sym(M)$, aljabar $\Sym(M)$ adalah $R$-aljabar komutatif. Jadi aljabar simetris dapat dipandang sebagai komutativisasi $T(M)$.
	\item Perkalian dalam aljabar eksterior lazim ditulis $(x,y) \mapsto x \wedge y$. Perhatikan bahwa untuk setiap $x, y \in M$,
		\[ (x+y) \otimes (x+y) - x \otimes x - y \otimes y = x \otimes y + y \otimes x \in I_{\bigwedge}(M), \]
		rumus di atas mengakibatkan $x \wedge y = - y \wedge x$, yakni kasus antikomutativitas dalam Definisi \ref{def:anti-commutation} untuk unsur homogen berderajat satu. Karena $M$ membangun $\bigwedge(M)$, untuk setiap unsur homogen $x,y$ kita memperoleh $x \wedge y = (-1)^{\deg x \deg y} y \wedge x$. Sebaliknya, jika $2 \in R^\times$, antikomutativitas mengakibatkan $x \wedge x = - x \wedge x = 0$ untuk setiap $x \in M$; dalam hal ini aljabar eksterior dapat dipandang sebagai antikomutativisasi $T(M)$.
	\item Terakhir, untuk $\bigwedge(M)$ kita verifikasi \eqref{eqn:strong-anti-commutation}. Kasus derajat satu sudah jelas, sedangkan setiap unsur homogen berderajat ganjil dapat ditulis $\omega = \sum_{i=1}^k t_i \wedge \eta_i$, dengan $\deg(t_i)=1$ dan $\deg(\eta_i) \in 2\Z$. Antikomutativitas memberikan
		\begin{align*}
			\omega \wedge \omega & = \sum_{i,j} t_i \wedge \eta_i \wedge t_j \wedge \eta_j = \sum_{i,j} t_i \wedge t_j \wedge \eta_i \wedge \eta_j \\
			& = \sum_i (t_i \wedge t_i) \wedge (\eta_i \wedge \eta_i) + \sum_{i < j} (t_i \wedge t_j + t_j \wedge t_i) \wedge \eta_i \wedge \eta_j \\
			& = 0 + 0 = 0.
		\end{align*}
\end{itemize}
Hasil-hasil tersebut dirangkum sebagai berikut.
\begin{lemma}\label{prop:Sym-wedge-properties}
	Aljabar simetris $\Sym(M)$ adalah aljabar komutatif bergradasi; aljabar eksterior $\bigwedge(M)$ adalah aljabar antikomutatif bergradasi dan memenuhi sifat \eqref{eqn:strong-anti-commutation}.
\end{lemma}

\begin{example}\label{eg:Sym-wedge-rank1}
	Misalkan $M$ modul bebas ber-rank satu, $M = Rx$. Maka, sebagai aljabar bergradasi, $\Sym(M) = R[x] = \bigoplus_{n \geq 0} R x^n$, sedangkan $\bigwedge(M) = R \oplus Rx$.
\end{example}

\begin{lemma}
	Konstruksi aljabar simetris dan aljabar eksterior komutatif dengan perubahan basis: untuk homomorfisme gelanggang $R \to S$, isomorfisme fungtor dalam Lemma \ref{prop:tensor-alg-basechange}, yaitu $T(- \otimes S) \rightiso T(-) \otimes S$, secara alami menginduksi
	\[ \Sym(- \otimes S) \rightiso \Sym(-) \otimes S, \quad \bigwedge(- \otimes S) \rightiso \bigwedge(-) \otimes S, \]
	dan isomorfisme-isomorfisme ini mempertahankan struktur bergradasi.
\end{lemma}
\begin{proof}
	Dengan menggabungkan Lemma \ref{prop:tensor-alg-basechange} dan deskripsi para pembangun ideal, langsung diperoleh bahwa $T(- \otimes S) \rightiso T(-) \otimes S$ menginduksi isomorfisme bergradasi $I_{\Sym}(M \otimes S) \rightiso I_{\Sym}(M) \otimes S$ dan $I_{\bigwedge}(M \otimes S) \rightiso I_{\bigwedge}(M) \otimes S$. Hasil kali tensor mempertahankan hasil bagi (Proposisi \ref{prop:tensor-mod-exact}), sehingga bukti selesai.
\end{proof}

Seperti hasil kali tensor multipeubah, pangkat simetris ke-$n$ dari $M$, $\Sym^n(M)$, dan pangkat eksterior $\bigwedge^n(M)$ juga dicirikan oleh sifat universal. Misalkan $A$ adalah sebarang $R$-modul dan $n \geq 1$; tetapkan
\begin{align*}
	\text{Sym}(M \times n; A) & := \left\{ B \in \text{Mul}(\underbracket{M, \ldots, M}_{n \text{ salinan}}; A) : B(\ldots, x,y, \ldots) = B(\ldots, y,x, \ldots) \right\}, \\
	\text{Alt}(M \times n; A) & := \left\{ B \in \text{Mul}(\underbracket{M, \ldots, M}_{n \text{ salinan}}; A) : B(\ldots, x,x, \ldots)=0 \right\}
\end{align*}
Dalam rumus definisi ini, $x,y \in M$, dan kedua posisi bersebelahan yang ditampilkan boleh sebarang. Unsur-unsur $\text{Sym}(M \times n; A)$ dan $\text{Alt}(M \times n; A)$ masing-masing disebut pemetaan multilinear simetris dan selang-seling berderajat $n$ pada $M$ dengan nilai di $A$. Untuk menghubungkannya dengan konsep serupa dalam aljabar linear, perhatikan dahulu grup simetris $\mathfrak{S}_n$: melalui $(\sigma B)(x_1, \ldots, x_n) = B(x_{\sigma(1)}, \ldots, x_{\sigma(n)})$, grup ini beraksi dari kiri pada $\text{Mul}(M, \ldots, M; A)$. Ingat pula bahwa $\mathfrak{S}_n$ dibangun oleh transposisi $\tau_i := (i \quad i+1)$, $1 \leq i < n$ (Lemma \ref{prop:S_n-generation}).
\begin{compactitem}
	\item Untuk $\text{Sym}(M \times n; A)$, definisi dapat ditulis ulang sebagai $\forall i, \;\tau_i B = B$, atau lebih lanjut sebagai $\forall \sigma \in \mathfrak{S}_n,\; \sigma B =B$. Inilah definisi yang wajar bagi pemetaan multilinear simetris.
	\item Untuk $\text{Alt}(M \times n; A)$, dengan meniru penurunan antikomutativitas $\bigwedge(M)$ di atas, dari
	\[ B(\ldots, x+y, x+y, \ldots, \ldots) = B(\ldots, x,x, \ldots) = B(\ldots, y,y, \ldots) = 0\]
	diperoleh $B(\ldots, x,y, \ldots) = -B(\ldots, y,x, \ldots)$, yakni $\tau_i B = -B$; bahkan lebih lanjut:
		\begin{gather*}
			\forall \sigma \in \mathfrak{S}_n,\; \sigma B = \sgn(\sigma)B.
		\end{gather*}
		Jika $2 \in R^\times$, dengan mengambil $x=y$ kita juga dapat menyimpulkan kembali $B(\ldots, x,x, \ldots)=0$. Jadi $\text{Alt}(M \times n; -)$ pun merupakan definisi yang wajar bagi pemetaan multilinear selang-seling, setidaknya tanpa keraguan ketika $2 \in R^\times$.
\end{compactitem}
Untuk $n=0$, definisikan $\text{Sym}(M \times 0; A) = \text{Alt}(M \times 0; A) = A$; perhatikan bahwa $\Sym^0 = \bigwedge^0 = R$.

\begin{proposition}\label{prop:Symm-Alt-univ}
	Untuk setiap $M$ dan $n \geq 1$ seperti di atas, terdapat isomorfisme fungtor
	\begin{align*}
		\Hom_{R\dcate{Mod}}\left( \Sym^n(M), - \right) & \longrightiso \mathrm{Sym}(M \times n; -) \\
		\varphi & \longmapsto [(x_1, \ldots, x_n) \mapsto \varphi(x_1 \cdots x_n)], \\
		\Hom_{R\dcate{Mod}}(\bigwedge^n(M), - ) & \longrightiso \mathrm{Alt}(M \times n; -) \\
		\varphi & \longmapsto [(x_1, \ldots, x_n) \mapsto \varphi(x_1 \wedge \cdots \wedge x_n)].
	\end{align*}
	Untuk $n=0$, isomorfisme di atas tetap berlaku berdasarkan definisi.
\end{proposition}
\begin{proof}
	Abaikan kasus trivial $n=0$. Untuk kasus $\Sym^n(M)$, perhatikan bahwa $I^n_{\Sym}(M)$ dibangun oleh unsur-unsur berbentuk
	\[ a(x \otimes y )b- a(y \otimes x)b, \quad x,y \in M, \; a,b \in T(M): \text{unsur homogen}, \quad \deg(a)+\deg(b)+2=n. \]
	Dalam sifat universal hasil kali tensor multipel \eqref{eqn:Mul-tensor}, $\varphi: T^n(M) \to A$ bernilai nol pada $I^n_{\Sym}(M)$ jika dan hanya jika $B \in \text{Mul}(M, \cdots, M; A)$ yang bersesuaian memenuhi
	\[ B(\ldots, x, y, \ldots) = \varphi(\cdots \otimes x \otimes y \otimes \cdots) = \varphi(\cdots \otimes y \otimes x \otimes \cdots) = B(\ldots, y, x, \ldots). \]
	Dari sini isomorfisme yang diinginkan segera tampak. Kasus selang-seling dapat diuraikan dengan cara yang sama.
\end{proof}

Perkenalkan notasi $R\dcate{CAlg}_\Z$ untuk kategori aljabar-$R$ komutatif bergradasi $\Z$, dan $R\dcate{CAlg}^-_\Z$ untuk kategori aljabar-$R$ antikomutatif yang memenuhi \eqref{eqn:strong-anti-commutation} dan bergradasi $\Z$. Kita segera memperoleh hasil berikut.
\begin{theorem}\label{prop:Sym-wedge-universal}
	Homomorfisme modul $M \to \Sym(M)$ dan $M \to \bigwedge(M)$ menginduksi isomorfisme fungtor
	\begin{align*}
		\Hom_{R\dcate{CAlg}_\Z}(\Sym(-),-) & \longrightiso \Hom_{R\dcate{Mod}}(-, U(-)), \\
		\Hom_{R\dcate{CAlg}^-_\Z}(\bigwedge(-),-) & \longrightiso \Hom_{R\dcate{Mod}}(-, U(-)),
	\end{align*}
	dengan $U$ fungtor $R\dcate{CAlg}_\Z, R\dcate{CAlg}^-_\Z \to R\dcate{Mod}$ yang memetakan $A$ ke $A_1$.
\end{theorem}
\begin{proof}
	Ambil kasus aljabar simetris sebagai contoh. Isomorfisme yang ditampilkan memetakan $f: \Sym(M) \to A$ ke komposisi $M \hookrightarrow \Sym(M) \xrightarrow{f} A$. Arah sebaliknya memetakan homomorfisme $R$-modul $\varphi: M \to A_1$ ke $f: \Sym(M) \to A$, dengan $f(x_1 \cdots x_n) = \varphi(x_1) \cdots \varphi(x_n)$; bahwa pemetaan terakhir terdefinisi dengan baik merupakan akibat Proposisi \ref{prop:Symm-Alt-univ}.
\end{proof}

Proposisi \ref{prop:graded-alg-cocomplete} memastikan bahwa $R\dcate{CAlg}_\Z$ dan $R\dcate{CAlg}^-_\Z$ keduanya kategori kokomplet. Dengan Teorema \ref{prop:Sym-wedge-universal}, bukti pernyataan berikut sepenuhnya sama dengan bukti Korolari \ref{prop:tensor-alg-limit}.
\begin{corollary}\label{prop:Sym-wedge-rightexact}
	Funktor $\Sym: R\dcate{Mod} \to R\dcate{CAlg}_\Z$ mempertahankan $\varinjlim$, dan epimorfisme modul $M \to M/N$ dipetakan ke $\Sym(M) \twoheadrightarrow \Sym(M)/\lrangle{N} \simeq \Sym(M/N)$. Hal yang sama berlaku untuk fungtor $\bigwedge: R\dcate{Mod} \to R\dcate{CAlg}^-_\Z$.
\end{corollary}

Jumlah langsung dan hasil kali tensor masing-masing merupakan koproduk di dalam $R\dcate{Mod}$ dan di dalam $R\dcate{CAlg}_\Z$ atau $R\dcate{CAlg}^-_\Z$. Penjabaran definisi secara saksama memberikan hasil berikut.
\begin{corollary}\label{prop:Sym-wedge-direct-sum}
	Konstruksi aljabar simetris dan aljabar eksterior mengubah jumlah langsung menjadi hasil kali tensor: untuk setiap keluarga $R$-modul $\{M_i\}_{i \in I}$, tetapkan $M := \bigoplus_{i \in I} M_i$. Terdapat isomorfisme alami aljabar-$R$ bergradasi $\bigotimes_{i \in I} \Sym(M_i) \longrightiso \Sym(M)$ dan $\bigotimes_{i \in I} \bigwedge(M_i) \longrightiso \bigwedge(M)$, dengan $\bigotimes_i \bigwedge(M_i)$ dilengkapi struktur aljabar-$R$ bergradasi menurut kaidah tanda Koszul (Definisi--Teorema \ref{def:Koszul-sign-alg}). Keduanya dicirikan oleh diagram komutatif berikut: untuk setiap $j \in I$,
	\[\begin{tikzcd}[column sep=tiny, row sep=small]
		\bigotimes_{i \in I} \Sym(M_i) \arrow[rr, "\sim"] & & \Sym(M) \\
		& \Sym(M_j) \arrow[hookrightarrow, lu] \arrow[ru, "{\Sym(M_j \to M)}"'] &
	\end{tikzcd}\quad\begin{tikzcd}[column sep=tiny, row sep=small]
		\bigotimes_{i \in I} \bigwedge(M_i) \arrow[rr, "\sim"] & & \bigwedge(M) \\
		& \bigwedge(M_j) \arrow[hookrightarrow, lu] \arrow[ru, "{\bigwedge (M_j \to M)}"'] &
	\end{tikzcd}\]
\end{corollary}

\begin{corollary}\label{prop:Sym-wedge-free}
	Misalkan $M = \bigoplus_{x \in X} Rx$ adalah $R$-modul bebas dengan basis himpunan $X$. Maka
	\begin{compactitem}
		\item sebagai aljabar-$R$ bergradasi, $\Sym(M)$ isomorfik dengan aljabar polinomial $R[X]$;
		\item pilih sebarang urutan total pada $X$; sebagai $R$-modul, $\bigwedge(M)$ adalah modul bebas dengan basis
		\[ \left\{ x_1 \wedge \cdots \wedge x_k : \; k \geq 0, \; x_1 < \ldots < x_k \in X \right\} \]
		yang ditampilkan di atas.
	\end{compactitem}
	Jika $R$-modul $N$ dibangun oleh $n$ unsur, maka $m > n \implies \bigwedge^m N = \{0\}$.
\end{corollary}
\begin{proof}
	Pernyataan tentang modul bebas $M$ dapat direduksi melalui Korolari \ref{prop:Sym-wedge-direct-sum} ke kasus rank satu, yakni Contoh \ref{eg:Sym-wedge-rank1}. Misalkan $R$-modul $N$ dilengkapi epimorfisme $R^{\oplus n} \twoheadrightarrow N$. Korolari \ref{prop:Sym-wedge-rightexact} mengakibatkan homomorfisme bergradasi $\bigwedge(R^{\oplus n}) \twoheadrightarrow \bigwedge(N)$ surjektif, sehingga pernyataan tentang $\bigwedge^m N$ tereduksi ke kasus modul bebas.
\end{proof}

Terakhir, mari kita lihat kaitannya dengan definisi yang lazim dalam geometri diferensial dan bidang-bidang lain. Sifat universal dalam Proposisi \ref{prop:Symm-Alt-univ} menunjukkan bahwa mendefinisikan $\Sym(M)$ dan $\bigwedge(M)$ sebagai hasil bagi $T(M)$ memang wajar. Akan tetapi, dalam beberapa sumber (misalnya \cite[\S 2.2]{ChCh}), keduanya didefinisikan sebagai modul bagian $T(M)$ yang terdiri atas semua tensor simetris $T_{\Sym}(M)$ dan tensor antisimetris $T_\wedge(M)$, bukan sebagai modul hasil bagi. Kita ulas definisi-definisi itu secara singkat. Perhatikan bahwa grup simetris $\mathfrak{S}_n$ beraksi dari kiri pada $T^n(M)$ melalui $\sigma(x_1 \otimes \cdots \otimes x_m) = x_{\sigma^{-1}(1)} \otimes \cdots \otimes x_{\sigma^{-1}(m)}$. Untuk sebarang homomorfisme grup $\chi: \mathfrak{S}_n \to \{\pm 1\}$, definisikan
\begin{gather*}
	T^n_\chi(M) := \{x \in T^n(M): \forall \sigma \in \mathfrak{S}_n,\; \sigma x = \chi(\sigma)x\}, \\
	T_\chi(M) := \bigoplus_{n \geq 1} T^n_\chi(M), \\
	T_{\Sym}(M) := T_1(M), \quad T_\wedge := T_{\sgn}(M).
\end{gather*}
Pandang setiap $\sigma$ sebagai automorfisme modul pada $T^n(M)$. Jika $n! \in R^\times$, mudah dilihat bahwa unsur di dalam $\End_R(T^n(M))$
\[ e_\chi = e^n_\chi := \frac{1}{n!} \sum_{\sigma \in \mathfrak{S}_n} \chi(\sigma)^{-1} \sigma \]
memenuhi
\begin{gather}\label{eqn:e_chi-equivariance}
	e_\chi|_{T^n_\chi(M)} = \identity, \qquad e_\chi \tau = \chi(\tau) e_\chi = \tau e_\chi, \quad \tau \in \mathfrak{S}_n.
\end{gather}
Dari sini diperoleh $\Image(e_\chi) = T^n_\chi(M)$, lalu $e_\chi^2 = e_\chi$. Menurut teori dalam \S\ref{sec:indecomposable-mod}, kita mendapatkan
\begin{align*}
	T^n(M) & = T^n_\chi(M) \oplus \Ker(e_\chi) \\
	x & = e_\chi(x) + (1 - e_\chi)(x).
\end{align*}
Selanjutnya boleh kita anggap $n \geq 2$. Kondisi $n! \in R^\times$ mengakibatkan $2 \in R^\times$; berdasarkan pembahasan pada awal bagian ini dan \eqref{eqn:e_chi-equivariance},
\begin{align*}
	I^n_{\Sym}(M) & = \sum_{\tau \in \mathfrak{S}_n} \Image(\tau - 1) \subset \Ker(e_1), \\
	I^n_{\bigwedge}(M) & = \sum_{\tau \in \mathfrak{S}_n} \Image(\tau - \sgn(\tau)) \subset \Ker(e_{\sgn});
\end{align*}
sebenarnya cukup mengambil $\tau= (i \quad i+1)$ sebagai pembangun; perhatikan bahwa rumus di atas trivial untuk $n=1$. Di pihak lain, citra $1 - e_\chi = \frac{1}{n!} \sum_\tau (1 - \chi(\tau)^{-1} \tau)$, yakni $\Ker(e_\chi)$, termuat di dalam $I^n_{\cdots}(M)$ (ambil $\chi = 1, \sgn$, sehingga $\chi=\chi^{-1}$). Dengan demikian kita sampai pada hasil berikut.
\begin{theorem}
	Jika $n! \in R^\times$, maka $\Ker(e_1) = I^n_{\Sym}$ dan $\Ker(e_{\sgn}) = I^n_{\bigwedge}$; dengan demikian identitas menginduksi isomorfisme $R$-modul
	\begin{gather*}
		T^n_{\Sym}(M) \longrightiso \Sym^n(M), \\
		T^n_\wedge(M) \longrightiso \bigwedge^n(M).
	\end{gather*}
\end{theorem}

Jika $R \supset \Q$, aljabar simetris dan aljabar eksterior dapat dibenamkan sebagai $R$-modul ke dalam $T(M)$. Pada unsur homogen, perkalian tensor simetris (atau antisimetris) diterjemahkan menjadi $(x,y) \mapsto e^{\deg x + \deg y}_1(x \otimes y)$ (atau $(x, y) \mapsto e^{\deg x + \deg y}_{\sgn}(x \otimes y)$), mungkin dengan selisih suatu faktor yang bergantung pada konvensi; lihat \cite[\S 2 (3.1)]{ChCh}. Dalam hal ini
\begin{gather*}
	T^2(M) = \Sym^2(M) \oplus \bigwedge^2(M), \\
	x \otimes y = x \cdot y + x \wedge y, \quad x,y \in M.
\end{gather*}

\section{Contoh penerapan: varietas Grassmann}\label{sec:Grassmannian}
Misalkan $F$ medan dan $V$ ruang vektor taknol berdimensi hingga atas $F$. Tetapkan $n := \dim_F V$. Tujuan bagian ini adalah mempelajari himpunan semua subruang vektor $W$ berdimensi $k$ di dalam $V$, yaitu $\Gras(k, V)$, dengan $0 \leq k \leq n$. Lebih tepatnya, kita hendak memperoleh parameterisasi yang mudah bagi unsur-unsur $\Gras(k,V)$ serta memberinya struktur geometris yang sesuai. Objek geometris ini disebut \emph{varietas Grassmann}; ia merupakan alat yang sangat penting dalam banyak bidang, misalnya: \index{Varietas Grassmann (Grassmannian)}\index[sym1]{G(k,V)@$\Gras(k,V)$}
\begin{compactenum}[(a)]
	\item Dalam mempelajari submanifold tertutup berdimensi $k$, $M$, dari $\R^n$ (misalnya permukaan mulus), ruang singgung di setiap titik membentuk pemetaan mulus dari $M$ ke $\Gras(k, \R^n)$ atau ke versinya yang berorientasi. Pemetaan ini disebut pemetaan Gauss dan memuat banyak informasi mengenai geometri submanifold tersebut.
	\item Dalam teori homotopi, grup uniter $\mathrm{U}(k)$ mempunyai ruang pengklasifikasi $\mathbf{B}\mathrm{U}(k)$ yang dapat dibangun sebagai limit induktif $\varinjlim_{n \geq k}\Gras(k, \CC^n)$.
\end{compactenum}
Aljabar eksterior dari Definisi \ref{def:Sym-wedge} menyediakan alat yang ampuh untuk menangani persoalan-persoalan ini. Pertama-tama kita telaah beberapa kasus khusus yang sederhana.
\begin{itemize}
	\item Abaikan kasus trivial $k=0,n$.
	\item Jika $k=1$, maka $\Gras(1, V)$ adalah himpunan semua garis di dalam $V$. Jadi terdapat bijeksi
		\begin{align*}
			(V \smallsetminus \{0\}) \big/ F^\times & \xrightarrow{1:1} \Gras(1, V) \\
			F^\times \cdot v & \mapsto Fv
		\end{align*}
		Objek yang diperoleh dengan cara ini disebut ruang proyektif $\PP(V)$. Kita mengetahui bahwa $\PP(V)$ memiliki struktur geometris yang baik: pilih suatu basis untuk mengidentifikasi $V$ dengan $F^{n+1}$, dan tuliskan $\PP^n := \PP(F^{n+1})$. Garis yang direntang oleh $(x_0, \ldots, x_n) \in F^{n+1} \smallsetminus \{0\}$ ditulis $(x_0: \cdots :x_n) \in \PP^n$. Maka $\PP^n = \bigcup_{i=0}^n U_i$, dengan $U_i := \{(x_0: \cdots :x_n) : x_i \neq 0 \}$. Setiap unsur di petak $U_i$ dapat ditulis secara tunggal sebagai $(x_0: \cdots: \underbracket{1}_{i}: \cdots x_n)$, sehingga terdapat bijeksi $\varphi_i: U_i \rightiso F^n$. Selanjutnya, untuk $i \neq i'$, mudah diperiksa pada $U_i \cap U_{i'}$ bahwa
		\[ \varphi_{i'} \circ \varphi_i^{-1}\left( (x_j)_{j \neq i} \right) = \left( \frac{x_j}{x_{i'}} \right)_{j \neq i'}. \]
		Jika $F=\R, \CC$, segera tampak bahwa atlas koordinat $(U_i, \varphi_i)_{i=0}^n$ memberi $\PP^n$ struktur manifold diferensiabel atau manifold kompleks. Untuk $F$ umum, dalam bahasa geometri aljabar dapat dikatakan bahwa $\PP^n$ membentuk varietas aljabar di atas $F$. Ruang-ruang ini merupakan objek dasar dalam geometri; pembahasan terperinci dapat ditemukan dalam \cite[\S 5.3]{Xi18}.
	\item Pembenaman ruang vektor $V \hookrightarrow V'$ secara alami menginduksi $\Gras(k, V) \hookrightarrow \Gras(k, V')$; sebagai kasus khusus, $\PP(V) \hookrightarrow \PP(V')$.
	\item Memberikan subruang $W$ berdimensi $k$ ekuivalen dengan memberikan subruang berdimensi $n-k$ dalam ruang dual $V^\vee = \Hom_F(V,F)$, yaitu $W^\perp := \{ \check{v} \in V^\vee : \check{v}|_W = 0 \}$. Karena itu terdapat bijeksi alami
		\[ \Gras(k, V) \simeq \Gras(n-k, V^\vee). \]
		Khususnya, $\Gras(n-1, V)$ dapat diidentifikasi dengan ruang proyektif $\PP(V^\vee)$.
\end{itemize}

\begin{definition-theorem}[Pembenaman Plücker] \index{Pembenaman Plücker}
	Untuk $V$ dan $k$ seperti di atas, rumus berikut memberikan injeksi yang terdefinisi dengan baik:
	\begin{align*}
		\psi: \Gras(k,V) & \longrightarrow \PP( \bigwedge^k V) \\
		W & \longmapsto F^\times \cdot w_1 \wedge \cdots \wedge w_k = \bigwedge^k W \smallsetminus \{0\}
	\end{align*}
	dengan $w_1, \ldots, w_k$ sebarang basis dari $W$.
\end{definition-theorem}
\begin{proof}
	Pertama kita tunjukkan bahwa $\psi$ terdefinisi dengan baik. Pilih sebarang subruang $U$ sehingga $V = U \oplus W$. Korolari \ref{prop:Sym-wedge-direct-sum} menyatakan bahwa untuk setiap $h$ terdapat isomorfisme linear
	\begin{align*}
		\bigoplus_{a+b=h} (\bigwedge^a U \otimes \bigwedge^b W) & \longrightiso \bigwedge^h V \\
		\eta \otimes \xi & \longmapsto \eta \wedge \xi.
	\end{align*}
	Ambil $h=k$. Korolari \ref{prop:Sym-wedge-free} memastikan bahwa $w_1 \wedge \cdots \wedge w_k$ merupakan unsur taknol dalam suku jumlah langsung $(a,b)=(0,k)$ dan merentang garis $\bigwedge^k W$. Tuliskan $F^\times \Lambda := \psi(W)$. Pengamatan di atas bahwa $\wedge: \bigwedge U \otimes \bigwedge W \rightiso \bigwedge V$ mengakibatkan
	\begin{gather}\label{eqn:Plucker-recovery}
		W = \left\{ v \in V : v \wedge \Lambda = 0 \right\}
	\end{gather}
	Dengan demikian, $W$ dapat dipulihkan dari $\psi(W)$.
\end{proof}

Pembenaman $\psi$ saja belum cukup untuk menampakkan sifat geometris $\Gras(k, V)$; kita ingin mencirikan citra $\psi$ setepat mungkin. Untuk itu diperlukan pemahaman yang lebih mendalam tentang aljabar eksterior. Ingat bahwa $V^\vee$ adalah ruang dual dari $V$. Tuliskan $\lrangle{\cdot, \cdot}: V^\vee \times V \to F$ untuk pasangan alami.
\begin{definition-theorem}\label{def:wedge-contraction}\index{suobing@Kontraksi (contraction)}
	Terdapat tepat satu pemetaan linear $\iota: V^\vee \to \End_F( \bigwedge V)$ yang memenuhi sifat-sifat berikut:
	\begin{gather*}
		\iota(\check{v})(v) = \lrangle{\check{v}, v}, \quad v \in V = \bigwedge^1 V, \\
		\iota(\check{v})(\xi \wedge \eta) = \iota(\check{v})(\xi) \wedge \eta + (-1)^{\deg \xi} \xi \wedge \iota(\check{v})(\eta),
	\end{gather*}
	dengan $\xi, \eta \in \bigwedge V$ unsur homogen dalam aljabar eksterior; akibatnya, untuk setiap $h$ berlaku $\iota(\check{v})(\bigwedge^{h+1} V) \subset \bigwedge^h(V)$.
\end{definition-theorem}
\begin{proof}
	Jelas bahwa $\iota$ semacam itu tunggal; ia harus memenuhi
	\[ \iota(\check{v})(v_0 \wedge \cdots \wedge v_h) = \sum_{i=0}^k (-1)^i v_0 \wedge \cdots \wedge \lrangle{\check{v}, v_i} \wedge \cdots \wedge v_h \]
	dengan $v_0, \ldots, v_h \in V$. Sebaliknya, untuk memeriksa bahwa rumus ini sungguh memberikan $\iota(\check{v}) \in \Hom_F( \bigwedge^{h+1} V, \bigwedge^h V)$ yang terdefinisi dengan baik, cukup ingat kembali konstruksi $\bigwedge V$ dalam Definisi \ref{def:Sym-wedge}. Perinciannya diserahkan kepada pembaca yang berminat.
\end{proof}
Untuk memahami $\iota$ secara konkret, pilih basis $v_1, \ldots, v_n$, dan biarkan $\check{v}_1, \ldots, \check{v}_n$ menjadi basis dual dari $V^\vee$. Dengan demikian $\bigwedge^k V$ mempunyai basis berbentuk $\{ v_{i_1} \wedge \cdots \wedge v_{i_k}: 1 \leq i_1 < \cdots < i_k \leq n \}$. Cukup tinjau kasus $\check{v} = \check{v}_1$; pencirian di atas segera memberikan
\begin{gather}\label{eqn:contraction-concrete}
	\iota(\check{v}_1)(v_{i_1} \wedge \cdots \wedge v_{i_k}) = \begin{cases}
		v_{i_2} \wedge \cdots \wedge v_{i_k}, & i_1 = 1 \\
		0, & i_1 > 1. 
\end{cases}\end{gather}
Karena alasan ini, $\iota$ juga disebut operasi kontraksi pada aljabar eksterior. Dengan menerapkan kontraksi berulang kali, untuk setiap $\check{v}_1, \ldots, \check{v}_r \in V^\vee$ kita memperoleh $\iota(\check{v}_1) \cdots \iota(\check{v}_r) \in \End_F(\bigwedge V)$. Mudah dilihat bahwa ini multilinear dalam $(\check{v}_1, \ldots, \check{v}_r)$; lebih lanjut dapat dibuktikan bahwa ia menginduksi
\begin{equation}\label{eqn:contraction-general}
	\bigwedge^r(V^\vee) \to \End_F(\bigwedge V).
\end{equation}	
Salah satu cara cepat menurunkan \eqref{eqn:contraction-general} adalah, tanpa mengurangi keumuman, mengambil $\check{v} = \check{v}_1$ seperti dalam \eqref{eqn:contraction-concrete}; dari sini langsung tampak bahwa $\iota(\check{v}) \iota(\check{v})=0$. Mengingat kembali konstruksi Definisi \ref{def:Sym-wedge}, hal ini sudah memadai.

Kembali ke pembenaman Plücker. Untuk sebarang $F^\times \Lambda \in \PP(\bigwedge^k V)$, mula-mula kita cari solusi hampiran terbaik $W$ bagi persoalan $\psi(W) \stackrel{?}{=} F^\times \Lambda$.
\begin{lemma}\label{prop:Plucker-W}
	Misalkan $1 \leq k \leq n$ dan $\Lambda \in \bigwedge^k V$ taknol. Definisikan $W \subset V$ sebagai subruang yang direntang oleh semua $\{ \iota(\Xi) \Lambda : \Xi \in \Lambda^{k-1}(V^\vee) \}$. Maka, untuk setiap subruang $W_0 \subset V$,
	\[ \Lambda \in \Image\left[ \bigwedge^k W_0 \to \bigwedge^k V \right] \iff W_0 \supset W. \]
	Khususnya, dengan mengambil $W_0=W$ diperoleh $\dim W \geq k$, dan $\dim W = k$ mengakibatkan $\psi(W) = F^\times\Lambda$.
\end{lemma}
\begin{proof}
	Untuk $W_0$ yang diberikan, pilih $U$ sehingga $V = U \oplus W_0$; dengan demikian $V^\vee =U^\vee \oplus W_0^\vee$. Kita kembali mempunyai penguraian
	\[ \wedge: \bigoplus_{a+b=k} (\bigwedge^a U \otimes \bigwedge^b W_0) \longrightiso \bigwedge^k V.\]
	Karena itu, $\Lambda \in \Image[\bigwedge^k W_0 \to \bigwedge^k V]$ jika dan hanya jika semua komponen $\Lambda$ dengan $a > 0$ bernilai nol. Pilih basis
	\[ \underbracket{u_1, \ldots, u_{n-r}}_{U\; \text{: basis}}, \underbracket{w_1, \ldots, w_r}_{W_0\; \text{: basis}}. \]
	Andaikan $\Lambda$ mempunyai komponen taknol untuk suatu $a > 0$, misalnya setelah dijabarkan terhadap basis, koefisien
	\[ u_{i_1} \wedge \cdots \wedge u_{i_a} \wedge w_{i_{a+1}} \wedge \cdots \wedge w_{i_k} \]
	taknol. Dengan menggunakan basis dual, $\Xi := \check{u}_{i_2} \wedge \cdots \wedge \check{w}_{i_k}$ mengontraksikan $\Lambda$ menjadi kelipatan taknol dari $u_{i_1}$ menurut \eqref{eqn:contraction-concrete}; jadi $W \not\subset W_0$. Sebaliknya, jika $\Lambda \in \bigwedge^k W_0$, maka kontraksi apa pun atasnya tetap berada di dalam $\bigwedge^{\leq k} W_0$, sehingga $W \subset W_0$. Sifat terakhir dalam pernyataan mengikuti dari $\dim W = k \implies \dim \bigwedge^k W = 1$.
\end{proof}

\begin{lemma}
	Untuk unsur taknol $\Lambda \in \bigwedge^k V$, definisikan $W$ seperti di atas dan $W' := \{ w \in W: w \wedge \Lambda = 0 \}$. Maka
	\[ F^\times \cdot \Lambda \in \Image(\psi) \iff W'=W; \]
	jika kondisi ini terpenuhi, $\psi(W) = F^\times \cdot\Lambda$.
\end{lemma}
\begin{proof}
	Jika $F^\times \cdot \Lambda = \psi(W_0)$, maka $\dim_F W_0 = k$ dan Lemma \ref{prop:Plucker-W} memberikan $W_0 \supset W$. Bersama $\dim_F W \geq k$, ini menghasilkan $W_0=W$; pengamatan \eqref{eqn:Plucker-recovery} segera memberikan $W'=W$. Sebaliknya, andaikan $W'=W$. Karena bentuk bilinear
	\[ \wedge: \bigwedge^{\dim W - k} W \otimes \bigwedge^k W \to \bigwedge^{\dim W} W \]
takdegenerat (terapkan Korolari \ref{prop:Sym-wedge-free}) dan $\Lambda \in \bigwedge^k W \smallsetminus \{0\}$, satu-satunya kemungkinan ialah $\dim W = k$; Lemma \ref{prop:Plucker-W} kemudian mengakibatkan $F^\times\Lambda = \psi(W)$.
\end{proof}

Kondisi $W'=W$ ekuivalen dengan $(\iota(\Xi)\Lambda) \wedge \Lambda = 0$ untuk semua $\Xi \in \bigwedge^{k-1} (V^\vee)$. Dengan demikian, kita memperoleh sistem persamaan pendefinisi bagi $\Gras(k,V) \hookrightarrow \PP(\bigwedge^k V)$ sebagai berikut.
\begin{theorem}[Relasi Plücker]
	Citra pembenaman Plücker $\psi: \Gras(k,V) \hookrightarrow \PP(\bigwedge^k V)$ sama dengan himpunan titik nol bersama dari keluarga fungsi kuadratik homogen
	\begin{align*}
		f_\Xi: \bigwedge^k V & \longrightarrow \bigwedge^{k+1} V \\
		\Lambda & \longmapsto (\iota(\Xi)\Lambda) \wedge \Lambda, \quad \Xi \in \bigwedge^{k-1} (V^\vee)
	\end{align*}
	yang ditampilkan di atas.
\end{theorem}
Karena homogen, himpunan titik nolnya invarian terhadap penskalaan dan dengan demikian mendefinisikan suatu subhimpunan dari $\PP(\bigwedge^k V)$. Setelah memilih basis dan menjabarkan nilai $f_\Xi$ menurut komponennya, sistem persamaan pendefinisi $\Image(\psi)$ dapat ditulis sebagai suatu keluarga polinomial kuadratik homogen. Dalam bahasa geometri aljabar, $\psi$ membenamkan $\Gras(k,V)$ sebagai varietas aljabar proyektif.

Sebagai contoh, ambil $k=2$ dan andaikan $2 \in F^\times$. Sifat kontraksi mengakibatkan $(\iota(\Xi) \Lambda) \wedge \Lambda = \frac{1}{2} \iota(\Xi)(\Lambda \wedge \Lambda)$; perhitungan kontraksi menggunakan \eqref{eqn:contraction-concrete} menunjukkan bahwa relasi Plücker menjadi $\Lambda \wedge \Lambda = 0$. Andaikan selanjutnya $\dim V = 4$; pilih basis $e_1, \ldots, e_4$ dan gunakan koordinat $\Lambda = \sum_{i < j} \lambda_{ij} e_i \wedge e_j$. Perhitungan langsung memberikan relasi Plücker
\[ \lambda_{12}\lambda_{34} - \lambda_{13}\lambda_{24} + \lambda_{14}\lambda_{23} = 0. \]
Dari sudut pandang geometri, ini menunjukkan bahwa $\Gras(2, V)$ dibenamkan sebagai hipersurfas kuadratik dalam $\PP(\bigwedge^2 V) \simeq \PP^5$; objek ini sangat menarik dalam geometri aljabar klasik.

\section{Determinan, teras, diskriminan}\label{sec:trace-norm-disc}
Misalkan $R$ gelanggang komutatif; selanjutnya tuliskan $\otimes := \otimes_R$.

Pertama, misalkan $E$ adalah $R$-modul bebas ber-rank $n$, dengan $n \in \Z_{\geq 0}$. Korolari \ref{prop:Sym-wedge-free} memastikan bahwa $\topwedge(E) := \bigwedge^n(E)$ adalah $R$-modul bebas ber-rank satu. Setiap endomorfismenya harus berbentuk perkalian skalar $x \mapsto rx$ (dengan $r \in R$ ditentukan secara tunggal), sehingga gelanggang $\End_R(\topwedge(E))$ dapat diidentifikasi dengan $R$. \index[sym1]{det(phi)@$\det \varphi$}
\begin{definition}\label{def:general-determinant} \index{hanglieshi@Determinan (determinant)}
	Misalkan $E$ adalah $R$-modul bebas ber-rank hingga. Untuk $\varphi \in \End_R(E)$, homomorfisme yang diinduksi oleh fungtorialitas ditulis $\det(\varphi) \in \End_R(\topwedge(E)) = R$.
\end{definition}
Begitu suatu basis $x_1, \ldots, x_n$ dari $E$ dipilih dan $\varphi(x_j) = \sum_{i=1}^n x_i a_{ij}$, Proposisi \ref{prop:Symm-Alt-univ} beserta sifat pemetaan multilinear selang-seling memberikan
\begin{multline}
	(\det\varphi)(x_1 \wedge \cdots \wedge x_n) = \varphi(x_1) \wedge \cdots \wedge \varphi(x_n) = \\
	\sum_{1 \leq \underbracket{j_1, \ldots, j_n}_{\text{saling berbeda}} \leq n} a_{j_1 1} a_{j_2 2} \cdots a_{j_n n} x_{j_1} \wedge \cdots \wedge x_{j_n} \\
	\left( \text{tetapkan}\; \sigma(i) = j_i \right) \qquad = \sum_{\sigma \in \mathfrak{S}_n} \sgn(\sigma) a_{\sigma(1) 1} \cdots a_{\sigma(n) n} x_1 \wedge \cdots \wedge x_n.
\end{multline}
Determinan matriks $A := (a_{ij})_{i,j}$ di atas $R$ didefinisikan sebagai unsur $R$
\begin{gather}\label{eqn:matrix-det}
	\det A := \sum_{\sigma \in \mathfrak{S}_n} \sgn(\sigma) a_{\sigma(1) 1} \cdots a_{\sigma(n) n},
\end{gather}
sehingga $\det \varphi$ teridentifikasi dengan $\det A$, sesuai dengan definisi yang lazim dalam aljabar linear. Sekarang mari kita lihat bagaimana sifat-sifat dasar determinan dapat diturunkan secara ringkas dari sudut pandang ini untuk sebarang gelanggang komutatif $R$.

\begin{theorem}\label{prop:matrix-det}
	Pilih basis $x_1, \ldots, x_n$ dari $R$-modul bebas $E$, dan gunakan hasil \S\ref{sec:free-modules} untuk mengidentifikasi unsur-unsur $\End_R(E)$ dengan unsur-unsur $M_n(R)$. Untuk setiap $A \in M_n(R)$, definisikan \emph{matriks adjugat}-nya sebagai
	\[ A^\vee := (A_{ji})_{\substack{1 \leq i \leq n \\ 1 \leq j \leq n}}, \quad A_{ij} := (-1)^{i+j} M_{ij} \]
	dengan $M_{ij} \in R$ determinan matriks yang diperoleh dari $A$ dengan menghapus baris ke-$i$ dan kolom ke-$j$; ia juga disebut minor.
	\begin{enumerate}[(i)]
		\item Untuk $\varphi, \psi \in \End_R(M)$ selalu berlaku $\det(\varphi\psi) = \det(\varphi) \det(\psi)$, dan $\det(\identity_E) = 1$.
		\item Transposisi matriks $A \mapsto {}^t A$ tidak mengubah determinan.
		\item Matriks adjugat memenuhi $A A^\vee = \det(A) \cdot 1_n = A^\vee A$.
		\item Endomorfisme $\varphi \in \End_R(E)$ invertibel jika dan hanya jika $\det \varphi \in R^\times$; dalam hal ini matriks $A$ yang bersesuaian memenuhi $A^{-1} = (\det A)^{-1} A^\vee$.
	\end{enumerate}
\end{theorem}
\begin{proof}
	(i) mengikuti dari fungtorialitas $\bigwedge^n(\cdot)$. (ii) merupakan akibat langsung dari \eqref{eqn:matrix-det} dan $\sgn(\sigma) = \sgn(\sigma^{-1})$. Pokok persoalannya adalah (iii). Karena transposisi membalik urutan perkalian, berdasarkan (ii) dan identitas mudah $({}^t A)^\vee = {}^t(A^\vee)$, cukup dibuktikan bahwa $A^\vee A = \det(A) \cdot 1_n$.

	Misalkan $A = (a_{ij})_{i,j}$ bersesuaian dengan $\varphi \in M_n(R)$. Di dalam $\bigwedge^n E$, jabarkan
	\[ \det\varphi(x_1 \wedge \cdots \wedge x_n) = \sum_{i=1}^n x_i a_{i1} \wedge \left[ \varphi(x_2) \wedge \cdots \wedge \varphi(x_n) \right]. \]
	Perhatikan bahwa $x_i \wedge \cdots \wedge x_i = 0$. Karena itu, setelah ruas kanan $[\cdots]$ dijabarkan, bagian yang memuat $x_i$ tidak menyumbang apa pun; perhitungannya tereduksi menjadi mencari determinan dari
	\[ R^{\oplus (n-1)} \xhookrightarrow{\text{komponen $1$ bernilai nol}} \underbracket{R^{\oplus n}}_{= E} \xrightarrow{\varphi} R^{\oplus n} \xrightarrow{\text{hapus komponen $i$}} R^{\oplus (n-1)} \]
	yang tepat sama dengan $M_{i1}$. Jadi
	\[ \det\varphi(x_1 \wedge \cdots \wedge x_n) = \sum_{i=1}^n a_{i1} x_i \wedge \left( M_{i1}  x_1 \wedge \cdots \widehat{x_i} \cdots \wedge x_n \right) \]
	dengan $\widehat{x_i}$ menyatakan bahwa suku tersebut dihapus. Setelah mengembalikan $x_i$ ke posisinya, diperoleh $\sum_i (-1)^{i+1} M_{i1} a_{i1} x_1 \wedge \cdots \wedge x_n$. Ini menunjukkan bahwa entri $(1,1)$ dari $A^\vee A$ sama dengan $\det A$. Dengan cara yang sama, setiap entri diagonal $A^\vee A$ sama dengan $\det A$.

	Selanjutnya kita buktikan bahwa entri ke-$(i,j)$ dari $A^\vee A$, yaitu $ \sum_k (-1)^{i+k} M_{ki} a_{kj}$, bernilai nol jika $i \neq j$. Perhatikan bahwa entri ini tidak bergantung pada kolom ke-$i$ dari $A$. Maka boleh kita andaikan kolom-kolom $i,j$ dari $A$ sama, yakni $a_{ki}=a_{kj}$; dengan demikian ekspresi di atas juga sama dengan entri ke-$(i,i)$ dari $A^\vee A$, yaitu $\det A$. Akan tetapi, citra endomorfisme $\varphi$ yang bersesuaian dengan $A$ termuat di dalam submodul $E' \subset E$ yang dibangun oleh $n-1$ unsur. Karena itu terdapat faktorisasi
	$\begin{tikzcd}[column sep=small] \bigwedge^n E \arrow[r] \arrow[rr, bend right="15", "\bigwedge^n \varphi" description] & \bigwedge^n E' \arrow[r] & \bigwedge^n E \end{tikzcd}$,
	Korolari \ref{prop:Sym-wedge-free} mengakibatkan suku tengah bernilai nol, sehingga $\det\varphi = 0$.

	Terakhir, kita buktikan (iv). Jika $\varphi$ invertibel, maka (i) mengakibatkan $\det\varphi \in R^\times$. Jika $\det\varphi \in R^\times$, maka (iii) secara eksplisit memberikan invers matriks $A$ yang bersesuaian.
\end{proof}

Untuk sebarang $R$-modul $E$, tetapkan $E^\vee := \Hom_R(E, R)$. Terdapat homomorfisme $R$-modul yang terdefinisi dengan baik
\begin{equation}\label{eqn:End-free-otimes}\begin{tikzcd}[row sep=tiny]
	R & \arrow[l] E^\vee \otimes E \arrow[r] & \End_R(E) \\
	f(x) & f \otimes x \arrow[mapsto, r] \arrow[mapsto, l] & \left[ v \mapsto f(v) x \right].
\end{tikzcd}\end{equation}

Misalkan $E$ modul bebas dan $n := \rank_R(E)$ berhingga. Lemma \ref{prop:Hom-matrix} dan pembahasan matriks di atas menunjukkan bahwa $E^\vee$ juga modul bebas ber-rank $n$, dan dalam \eqref{eqn:End-free-otimes} berlaku $E^\vee \otimes E \rightiso \End_R(E)$. Memang, setelah basis $x_1, \ldots, x_n$ dari $E$ dipilih dan didefinisikan
\[ \check{x}_i: a_1 x_1 + \cdots + a_n x_n \longmapsto a_i, \quad 1 \leq i \leq n \]
maka $\{\check{x}_1, \ldots, \check{x}_n\}$ adalah basis dari $E^\vee$ (disebut basis dual), dan diperoleh korespondensi tiga arah:
\[\begin{tikzcd}[row sep=small, column sep=tiny]
	\sum_{i,j} a_{ij} \check{x}_j \otimes x_i \arrow[phantom, d, sloped, "\in" description] & & \varphi: x_j \mapsto \sum_i x_i a_{ij} \arrow[phantom, d, sloped, "\in" description] \\
	E^\vee \otimes E \arrow[leftrightarrow, rr]  & & \End_R(E) \\
	& M_n(R) \arrow[leftrightarrow,lu] \arrow[leftrightarrow, ru] & \\
	& A = (a_{i,j})_{i,j} \arrow[phantom, u, sloped, "\in" description] &
\end{tikzcd} \]

Serupa dengan kasus ruang vektor, objek-objek berikut dapat didefinisikan bagi $\varphi \in \End_R(E)$ tanpa memilih basis:
\begin{itemize}
	\item Determinan $\det(\varphi) \in R$: telah dibahas dalam Definisi \ref{def:general-determinant};
	\item Teras $\Tr(\varphi) \in R$: ia tidak lain adalah citra $\varphi$ menurut \eqref{eqn:End-free-otimes} di bawah homomorfisme modul $\End_R(E) \rightiso E^\vee \otimes E \to R$; dari sudut pandang matriks, ini adalah rumus biasa $\Tr(A) = \sum_{i=1}^n a_{ii}$. \index[sym1]{Tr@$\Tr$}\index{ji-trace@Teras (trace)}
	\item Polinomial karakteristik $\mathrm{char}(\varphi, X) := \det(X \cdot \identity - \varphi) \in R[X]$, dengan $X$ sebagai peubah; di sini $\varphi$ diidentifikasi dengan unsur yang bersesuaian di dalam $\End_{R[X]}(E \otimes_R R[X])$.
\end{itemize}
Jika peran $E$ atau $R$ perlu ditonjolkan, tuliskan $\det_R(\varphi|E)$, $\Tr_R(\varphi|E)$, $\mathrm{char}_R(\varphi|E)$. Untuk setiap $\varphi, \psi \in \End_R(E)$, berlaku
\begin{align*}
	\det(\varphi\psi) & = \det(\varphi)\det(\psi), \quad \det(\identity_E)=1, \quad \det(0)=0 \\
	\det(r\varphi) &= r^{\rank_R(E)} \det(\varphi), \quad r \in R, \\
	\Tr(r\varphi + s\psi) & = r\cdot\Tr(\varphi) + s\cdot\Tr(\psi), \quad r,s \in R, \\
	\Tr(1) & = \rank_R(E). 
\end{align*}

Sekarang tinjau sebuah $R$-aljabar $A$, dan andaikan $A$, sebagai $R$-modul, bebas ber-rank hingga $n$.
\begin{definition}\label{def:norm-trace} \index{fanshu@norma (norm)} \index[sym1]{Tr}\index[sym1]{Nm@$\Nm$}
	Untuk $A$ seperti di atas dan $a \in A$, definisikan endomorfisme $R$-modul $m_a \in \End_R(A)$ sebagai perkalian kiri $x \mapsto ax$. Selanjutnya definisikan \emph{teras} $\Tr_{A|R}(a)$, \emph{norma} $\Nm_{A|R}(a)$, dan \emph{polinomial karakteristik} $\mathrm{char}_{A|R}(a, X)$ melalui
	\begin{align*}
		\Tr_{A|R}(a) & := \Tr(m_a), \\
		\Nm_{A|R}(a) & := \det(m_a), \\
		\mathrm{char}_{A|R}(a, X) & := \mathrm{char}(m_a, X) = \det_{R[X]}\left( m_{X-a} \mid A[X] \right) \\
		& = \Nm_{A[X]|R[X]}(X-a). \quad (\because\; A[X] \simeq A \otimes_R R[X])
\end{align*}
\end{definition}
Karena $a \mapsto m_a$ merupakan homomorfisme aljabar $A \to \End_R(A)$, determinan (atau teras) merupakan homomorfisme dari monoid multiplikatif (atau grup aditif) $A$ ke $R$. Pokok bagian ini ialah menetapkan transitivitas determinan dan teras.

\begin{lemma}\label{prop:free-transitivity}
	Misalkan $A$ suatu $R$-aljabar komutatif dan $E$ modul bebas atas $A$. Jika $A$, dipandang sebagai $R$-modul, adalah bebas, maka $E$, dipandang sebagai $R$-modul, juga bebas. Pilih basis $(e_i)_{i \in I}$ dari $E$ atas $A$ dan basis $(a_j)_{j \in J}$ dari $A$ atas $R$. Maka $(a_j e_i)_{(i,j) \in I \times J}$ membentuk basis dari $E$ atas $R$; akibatnya, $\rank_R(E) = \rank_R(A) \rank_A(E)$.
\end{lemma}
\begin{proof}
	Setiap $e \in E$ dapat ditulis secara tunggal sebagai kombinasi linear-$A$, $\sum_i u_i e_i$, dan setiap $u_i$ dapat ditulis secara tunggal sebagai kombinasi linear-$R$, $\sum_j r_{ij} a_j$. Karena itu, $(a_j e_i)_{i,j}$ benar-benar membentuk basis $R$-modul $E$: setiap $e$ mempunyai penyajian tunggal $e = \sum_{i,j} r_{ij} a_j e_i$.
\end{proof}

\begin{theorem}\label{prop:norm-trace-transitivity}
	Misalkan $A$ suatu $R$-aljabar komutatif yang sekaligus merupakan $R$-modul bebas ber-rank hingga, dan $E$ suatu $A$-modul bebas ber-rank hingga. Pandang $\varphi \in \End_A(E)$ yang diberikan sebagai unsur $\End_R(E)$. Maka
	\begin{gather*}
		\Tr_R(\varphi) = \Tr_{A|R}(\Tr_A(\varphi)), \quad \det_R(\varphi) = \Nm_{A|R}(\det_A(\varphi)), \\
		\mathrm{char}(\varphi, X) = \Nm_{A[X]|R[X]}(\mathrm{char}_A(\varphi, X)).
	\end{gather*}
\end{theorem}
\begin{proof}
	Pilih basis-$A$ $e_1, \ldots, e_n$ dari $E$ dan basis-$R$ $a_1, \ldots,  a_m$ dari $A$. Pertama, identifikasikan $\varphi$ melalui $\varphi(e_i) = \sum_{k=1}^n c_{ik} e_k$ dengan matriks $(c_{ik})_{1 \leq i,k \leq n} \in M_n(A)$. Maka
	\[ \varphi(a_j e_i) = a_j \varphi(e_i) = \sum_{k=1}^n a_j c_{ik} e_k. \]
	Selanjutnya $a_j c_{ik} = m_{c_{ik}}(a_j)$ dapat dijabarkan terhadap basis-$R$ $a_1, \ldots, a_m$. Ketika $j$ bervariasi, koefisien-koefisiennya tepat membentuk matriks persegi berorde $m$ yang bersesuaian dengan $m_{c_{ik}} \in \End_R(A)$. Jika $\varphi$ dinyatakan sebagai matriks terhadap basis-$R$ $\{ a_j e_i \}_{i,j}$, hasilnya dapat dipandang sebagai matriks persegi berorde $nm$ atas $R$ yang tersusun dari blok-blok $m \times m$, dengan blok ke-$(i,k)$ persis matriks yang bersesuaian dengan $m_{c_{ik}}$. Karena $A$ komutatif, kesemua $n^2$ blok ini saling komutatif terhadap perkalian. Persamaan untuk teras dan norma pun tereduksi ke Lemma aljabar linear \ref{prop:det-blocks}, sedangkan pernyataan tentang $\mathrm{char}(\varphi, X)$ tereduksi ke kasus norma.
\end{proof}

\begin{corollary}\label{prop:norm-trace-transitivity-2}
	Misalkan $B$ adalah $A$-aljabar yang sekaligus merupakan $A$-modul bebas ber-rank hingga. Maka setiap $b \in B$ memenuhi
	\begin{gather*}
		\Tr_{B|R}(b) = \Tr_{A|R}(\Tr_{B|A}(b)), \quad \Nm_{B|R}(b) = \Nm_{A|R}(\Nm_{B|A}(b)), \\
		\mathrm{char}_{B|R}(b,X) = \Nm_{A[X]|R[X]}(\mathrm{char}_{B|A}(b,X)).
	\end{gather*}
\end{corollary}
\begin{proof}
	Cukup gunakan $B$ sebagai pengganti $E$ dalam Teorema \ref{prop:norm-trace-transitivity}.
\end{proof}

\begin{definition}\label{def:discriminant}
	Misalkan $A$ suatu $R$-aljabar yang, sebagai $R$-modul, mempunyai basis $x_1, \ldots, x_n$. Diskriminan yang bersesuaian didefinisikan sebagai
	\[ d(x_1, \ldots, x_n) := \det_R \left( \Tr_{A|R}(x_i x_j) \right)_{1 \leq i,j \leq n} \quad \in R. \]
\end{definition}
Bentuk bilinear-$R$ simetris $(x,y) \mapsto \Tr_{A|R}(xy)$ disebut \emph{bentuk teras} dari $A$. Dua basis $x_i$ dan $y_j$ dihubungkan oleh $T = (t_{ij})_{i,j} \in \{M \in M_n(R): \det(M) \in R^\times \}$ melalui $y_i = \sum_j t_{ij} x_j$. Mudah dilihat bahwa \index{ji-trace!Bentuk teras (trace form)}
\[ d(y_1, \ldots, y_n) = \det(T)^2 d(x_1, \ldots, x_n) \]
Karena itu $d_A := d(x_1, \ldots, x_n) \mod R^{\times 2}\; \in R/R^{\times 2}$ hanya bergantung pada aljabar $A$ dan disebut \emph{diskriminan} dari $A$; khususnya, ideal utama yang dibangun oleh $d_A$ di dalam $R$ terdefinisi dengan baik dan disebut ideal diskriminan. Bentuk teras dan diskriminan yang diturunkan darinya merupakan invarian yang berguna dari aljabar $A$. \index{panbieshi}

Sekarang kita lengkapi bukti Teorema \ref{prop:norm-trace-transitivity}.
\begin{lemma}\label{prop:det-blocks}
	Misalkan matriks persegi $X$ berorde $nk$ di atas gelanggang komutatif $R$ ditulis dalam blok-blok $k \times k$ sebagai
	\[ X = \begin{pmatrix}
		a_{11} & \cdots & a_{1n} \\
		\vdots & \ddots & \vdots \\
		a_{n1} & \cdots & a_{nn}
	\end{pmatrix}, \quad a_{ij} \in M_k(R) \]
	dengan setiap $a_{ij}$ berada dalam suatu subaljabar komutatif $A$ dari $M_k(R)$. Jika $X$ dipandang sebagai unsur $M_n(A)$, tuliskan teras dan determinan yang diperoleh sebagai $\Tr_A(X)$ dan $\det_A(X)$, yang di bawah ini juga dipandang sebagai unsur $M_k(R)$. Tuliskan teras dan determinan aljabar matriks atas $R$ sebagai $\Tr_R$ dan $\det_R$. Maka
	\begin{align*}
		\Tr_R(X) & = \Tr_R( \Tr_A(X)), \\
		\det_R(X) & = \det_R(\det_A(X)).
	\end{align*}
\end{lemma}
\begin{proof}[{Bourbaki {\cite[III.112]{Bou-Alg1}}}]
	Kasus teras sudah sangat jelas. Untuk determinan, kita gunakan trik pelubangan matriks klasik. Kasus $n=1$ jelas dengan sendirinya, jadi selanjutnya andaikan $n \geq 2$. Tambahkan peubah formal $Z$ pada $R$ (yang dapat dipandang sebagai perturbasi terhadap $X$), dan nyatakan $X+Z$ sebagai matriks $n \times n$, $(\nu_{ij})_{i,j}$, atas $A[Z]$. Karena $A$ komutatif, matriks adjugat dari $X+Z$, yaitu $(X+Z)^\vee = (\check{\nu}_{ij})_{1 \leq i,j \leq n}$, dapat dibentuk di dalam $M_n(A[Z])$. Tujuannya adalah memperoleh persamaan matriks berikut di atas $A[Z]$:
	\begin{gather*} U := \begin{pmatrix}
		\check{\nu}_{11} & 0 & \cdots & 0 \\
		\check{\nu}_{12} & 1 & \cdots & 0 \\
		\vdots & 0 & 1 & \cdots \\
		\check{\nu}_{1n} & 0 & \cdots & 1
	\end{pmatrix}, \quad Q := \begin{pmatrix}
		\nu_{22} & \cdots & \nu_{2n} \\
		\vdots & \ddots & \\
		\nu_{n2} & \cdots & \nu_{nn}
	\end{pmatrix}\end{gather*}
	yang memenuhi
	\begin{gather*}
	(X+Z) U  = \begin{pmatrix}
		\det_{A[Z]}(X+Z) & \nu_{12} & \cdots & \nu_{1n} \\
		0 & \nu_{22} & \cdots & \nu_{2n} \\
		\vdots & \vdots & \ddots & \\
		0 & \nu_{n2} & \cdots & \nu_{nn}
	\end{pmatrix} = \begin{pmatrix}
		\det_{A[Z]}(X+Z) & \ast \\
		0 & Q
	\end{pmatrix}; \end{gather*}
	Jika dipandang sebagai matriks blok atas $R[Z]$, perkalian matriks blok menunjukkan bahwa persamaan di atas tetap berlaku. Kaidah perhitungan determinan segera memberikan
	\begin{gather*}
		\det_{R[Z]}(X+Z) \cdot \det_{R[Z]}(U) = \det_{R[Z]}\left( \det_{A[Z]}(X+Z)\right) \cdot \det_{R[Z]}(Q), \\
		\det_{R[Z]}(U) = \det_{R[Z]}(\check{\nu}_{11})
	\end{gather*}
	Secara rekursif, andaikan $\det_{R[Z]}(Q) = \det_{R[Z]}\left(\det_{A[Z]}(Q)\right)$. Definisi matriks adjugat atas $A[Z]$ memberikan $\det_{A[Z]}(Q) = \check{\nu}_{11}$. Jika $\det_{R[Z]}(\check{\nu}_{11})$ dapat dicoret dari kedua ruas persamaan pertama di atas, maka
	\[ \det_{R[Z]}(X+Z) = \det_{R[Z]}\left( \det_{A[Z]}(X+Z)\right), \]
	dan substitusi $Z=0$ memberikan hasil yang diinginkan, $\det_R(X) = \det_R(\det_A(X))$. Namun, dari definisi $\nu_{ij}$, $\det_{R[Z]}(\check{\nu}_{11}) = \det_{R[Z]}(Q)$ adalah polinomial monik dalam $Z$ berderajat $(n-1)k$, sehingga memang dapat dicoret. Terbuktilah pernyataan.
\end{proof}

% % %
\begin{Exercises}
	\item Misalkan $M$ modul atas gelanggang komutatif $R$. Definisikan $R$-modul $D(M) := R \oplus M$ dan lengkapi dengan perkalian $(r,m) (r', m') = (rr', rm' +r'm)$. Buktikan bahwa $D(M)$ menjadi $R$-aljabar.
	\item Berikan secara eksplisit isomorfisme $\CC$-aljabar $\CC \dotimes{\R} \CC \rightiso \CC \oplus \CC$.
	\item Verifikasikan bahwa objek yang didefinisikan dalam \eqref{eqn:quadratic-integer}, yaitu $\mathfrak{o}_D$, tepat merupakan himpunan bilangan bulat aljabar dalam $\Q(\sqrt{D})$.
	\item Misalkan $R$ daerah integral komutatif dan $K = \text{Frac}(R)$ medan pecahannya. Unsur $x \in K$ disebut hampir integral jika terdapat unsur taknol $u \in R$ sedemikian sehingga untuk setiap $n \geq 1$ berlaku $ux^n \in R$; jelas bahwa setiap unsur $R$ hampir integral. Buktikan bahwa
		\begin{compactitem}
			\item semua unsur hampir integral di dalam $K$ membentuk subgelanggang;
			\item jika $x \in K$ integral atas $R$, maka $x$ hampir integral; 
			\item jika $R$ gelanggang Noether (Definisi \ref{def:ACC-DCC-mod}), maka setiap unsur hampir integral adalah integral. \begin{hint} Dalam hal ini $R[x]$ adalah submodul dari $u^{-1} R$.\end{hint}
		\end{compactitem}
	\item Aljabar dalam soal ini boleh takasosiatif; ini berarti hukum asosiatif untuk perkalian dihilangkan dari Definisi \ref{def:algebra-mod-diagram}, tetapi keberadaan unsur satuan tetap disyaratkan. Tetapkan suatu medan $F$ dan tinjau $F$-aljabar berdimensi hingga $A$. Definisikan asosiator
		\[ (x,y,z) := (xy)z - x(yz), \quad x,y,z \in A. \]
		Jika selalu berlaku $(x,x,y)=0=(x,y,y)$, maka $A$ disebut \emph{aljabar alternatif}. Jika $\pi \in \End_F(A)$ memenuhi $\pi(xy)=\pi(y)\pi(x)$, $\pi(1_A)=1_A$, dan $\pi^2 = \identity_A$, maka objek ini disebut \emph{involusi}.
		\begin{enumerate}[(i)]
			\item Buktikan bahwa asosiator linear dalam setiap peubah; jika $A$ aljabar alternatif, buktikan $(x,y,z) = -(y,x,z)$ dan $(x,y,z) = -(x,z,y)$.
			\item Misalkan involusi $x \mapsto \bar{x}$ memenuhi $t(x) := x+\bar{x} \in F \cdot 1_A$ dan $n(x) := \bar{x}x = x\bar{x} \in F \cdot 1_A$. Maka setiap $x \in A$ memenuhi persamaan kuadrat $x^2 - t(x)x + n(x) 1_A = 0$.
			\item Misalkan $A$ alternatif dan tinjau involusi seperti di atas. Buktikan $(xy)\bar{y} = n(y)x$, $\bar{x}(xy) = n(x)y$, serta $n(xy)=n(x)n(y)$.
			\begin{hint} Untuk pernyataan terakhir, tulis ulang $n(xy)$ sebagai $((xy)\bar{y})\bar{x} - (xy, \bar{y}, \bar{x}) = n(x)n(y) - (\bar{x}, xy, \bar{y})$; persoalannya tereduksi menjadi membuktikan $(\bar{x}, xy, \bar{y}) = 0$. \end{hint}
			\item Misalkan $A$ suatu $F$-aljabar dan $n: A \to F$ pemetaan kuadratik yang memenuhi $n(xy)=n(x)n(y)$. Artinya, $n$ memenuhi $n(tx)=t^2 n(x)$ untuk $t \in F$, dan $B(x,y) := n(x+y)-n(x)-n(y)$ adalah bentuk bilinear simetris. Buktikan
				\begin{align*}
					B(xy,xy') = n(x)B(y,y') = B(yx,y'x), \\
					B(xy', x'y) + B(xy, x'y') = B(x,x')B(y,y');
				\end{align*}
				di bawah hipotesis butir sebelumnya, buktikan $B(xy,z) = B(y,\bar{x}z) = B(x, z\bar{y})$. \begin{hint} Dalam hal ini $t(x) = B(1,x)$ dan $\bar{x} = t(x) 1_A - x$; demikian pula untuk $y$.\end{hint}
		\end{enumerate}
	\item Dalam soal ini aljabar tetap boleh takasosiatif. Misalkan $A$ aljabar berdimensi hingga atas medan $F$, dan diberikan involusi $\pi: A \to A$ yang memenuhi
		\[ x + \pi(x) \in F \cdot 1_A, \quad x\pi(x) = \pi(x) x \in F \cdot 1_A \]
		Tetapkan $t(x), n(x)$ seperti dalam soal sebelumnya, dan andaikan $B_n(x,y) := n(x+y) - n(x) - n(y)$ takdegenerat:
		\[ B_n(x,-) = 0 \iff x = 0 \iff B_n(-,x)=0, \quad x \in A; \]
		Definisikan \emph{aljabar Cayley--Dickson} $\text{CD}(A, \lambda) := A \oplus vA$, dengan $v$ hanya sebuah simbol dan $\lambda \in F^{\times}$, serta perkalian
		\[ (a + vb) \cdot (a' + vb') = aa' + \lambda b' \pi(b) + v \left( \pi(a)b' + a'b \right). \]
		Definisikan $N(a+vb) = n(a) - \lambda n(b)$ dan $\Pi(a+vb) = \pi(a) - vb$.
		\begin{enumerate}[(i)]
			\item Buktikan bahwa $\text{CD}(A, \lambda)$ adalah $F$-aljabar dengan unsur satuan $1 := 1_A + v \cdot 0$, dan memuat $A$ sebagai subaljabar.
			\item Buktikan bahwa $\Pi$ adalah involusi pada $\text{CD}(A, \lambda)$, dan untuk setiap $x \in \text{CD}(A, \lambda)$ berlaku
				\begin{gather*}
					T(x) := x + \Pi(x) \in F \cdot 1, \quad N(x) = x\Pi(x) = \Pi(x) x \in F \cdot 1;
				\end{gather*}
			\item Kita mempunyai
				\begin{align*}
					\text{CD}(A, \lambda) \text{ alternatif} & \iff A \text { asosiatif}, \\
					\text{CD}(A, \lambda) \text{ asosiatif} & \iff A \text { asosiatif dan komutatif}, \\
					\text{CD}(A, \lambda) \text{ komutatif} & \iff A=F.
				\end{align*}
				\begin{hint} Untuk implikasi $\implies$ pertama, dari $N((a+vb)(c+vd)) = N(a+vb)N(c+vd)$ (hasil soal sebelumnya) dan $N(v)=-\lambda$, turunkan $B_n(ac, d\pi(b)) = B_n(\pi(a)d, cb) = B_n(cb, \pi(a)d)$, lalu peroleh $(ac)b = a(cb)$. Untuk implikasi $\implies$ kedua, perhatikan bahwa $v(ba) = (va)b$. \end{hint}
			\item Mulai dari $A=F=\R$ dan ambil $\lambda=-1$; secara bertahap kita memperoleh $\CC$, $\mathbb{H}$, dan suatu aljabar alternatif takasosiatif berdimensi $8$, $\mathbb{O}$. Selain itu, setiap $x \in \mathbb{O} \smallsetminus \{0\}$ mempunyai invers multiplikatif.
		\end{enumerate}
		Aljabar $\mathbb{O}$ lazim disebut aljabar oktonion. Sifat, penerapan, dan sejarah singkatnya dibahas terperinci dalam \cite{Baez02}. \index{bayuanshu}
	\item Untuk gelanggang komutatif $R$ dan ideal-idealnya $I, J$, buktikan $R/I \dotimes{R} R/J \simeq R/(I+J)$.
	\item Untuk sebarang medan $\Bbbk$, bangun isomorfisme alami
		\[ \Bbbk[ X_1^{\pm 1}, \ldots, X_n^{\pm 1}]^{\mathfrak{S}_n} \simeq \Bbbk[e_1, \ldots, e_{n-1}] \dotimes{\Bbbk} \Bbbk[e_n^{\pm 1}] \]
		dengan $e_1, \ldots, e_n$ polinomial simetris elementer dalam Definisi \ref{def:elementary-symm-poly}.
	\item Tetap misalkan $R$ gelanggang komutatif. Buktikan bahwa jika $M \simeq R/I_1 \oplus \cdots \oplus R/I_n$, dengan $I_1 \subset \cdots \subset I_n$ ideal-ideal, maka
		\[ \text{ann}\left( \bigwedge^a M \right) = I_a, \quad a=1, \ldots, n. \]
		Gunakan ini untuk menyederhanakan bukti ketunggalan dalam Teorema \ref{prop:elementary-divisor}, lalu perluas hasilnya ke sebarang gelanggang komutatif.
	\item Untuk ruang vektor berdimensi hingga $V$ atas medan $F$, definisikan $F[V]$ sebagai aljabar polinomial pada $V$: lebih tepatnya, setelah memilih basis $v_1, \ldots, v_n \in V$, $F[V]$ adalah aljabar polinomial dengan basis dual $\check{v}_1, \ldots, \check{v}_n \in V^\vee$ sebagai peubah; secara intrinsik, $F[V] = \Sym(V^\vee)$. Setiap $P \in F[V]$ dapat dievaluasi di suatu titik $v \in V$; nilainya ditulis $P_v \in F$. Selanjutnya andaikan $|F|$ takhingga dan $A$ suatu $F$-aljabar berdimensi hingga.
		\begin{compactenum}[(i)]
			\item Perkenalkan peubah $\lambda$ dan definisikan $\mathcal{P} := \left\{ P(\lambda) \in F[A][\lambda] : P \text{ monik}, \; \forall a \in A, \; P_a(a)=0 \right\}$. Buktikan bahwa terdapat tepat satu $P_A(\lambda) \in \mathcal{P}$ dengan $\deg_\lambda P_A(\lambda)$ minimal.
			\begin{hint}
				Untuk membuktikan bahwa $\mathcal{P}$ takkosong, tinjau $P_a(\lambda) := \det(\lambda \cdot \identity_A - L_a)$, dengan $L_a: A \to A$ aksi perkalian kiri dan simbol $a$ menyatakan unsur ``generik'' dari $A$; koefisien $P_a(\lambda)$ berada di dalam $F[A]$.
			\end{hint}
			\item Buktikan bahwa jika $Q(\lambda) \in F[A][\lambda]$ memenuhi $\forall a \; Q_a(a)=0$, maka $P_A(\lambda)$ membagi $Q(\lambda)$. Karena itu, $P_A(\lambda)$ disebut \emph{polinomial minimal generik} dari $A$.
			\item Untuk sebarang perluasan medan $E \supset F$, diperoleh $E$-aljabar $B := A_E$. Buktikan $P_A(\lambda) = P_B(\lambda)$. Dengan ini, perluas definisi $P_A(\lambda)$ ke kasus $F$ berhingga.
			\item Untuk $A = M_n(F)$ dan $A = \mathbb{H}$ (dalam kasus terakhir $F=\R$), tentukan $P_A(\lambda)$.
		\end{compactenum}
		Hasil di atas juga berlaku untuk aljabar berdimensi hingga, takasosiatif, dan berunsur satuan, asalkan setiap $a \in A$ membangun suatu aljabar asosiatif; ini mencakup kasus aljabar alternatif.
	\item Misalkan $M$ modul proyektif atas gelanggang komutatif $R$. Buktikan bahwa untuk setiap $n \geq 0$, modul $T^n M$, $\Sym^n M$, dan $\bigwedge^n M$ semuanya proyektif. \begin{hint}Proposisi \ref{prop:proj-mod-characterization} dapat digunakan.\end{hint}
	\item Buktikan Lemma Cartan berikut. Misalkan $V$ ruang vektor berdimensi hingga atas medan $F$, dan $v_1, \ldots, v_m \in V$ bebas linear. Buktikan bahwa jika $w_1, \ldots, w_m \in V$ memenuhi $\sum_{i=1}^m v_i \wedge w_i = 0$, maka terdapat keluarga $\left( a_{ij} \in F \right)_{1 \leq i,j \leq m}$ sedemikian sehingga $w_i = \sum_{j=1}^m a_{ij} v_j$ dan $a_{ij} = a_{ji}$.
	\item Dengan notasi soal sebelumnya, misalkan $\xi \in \bigwedge^k V$. Buktikan bahwa terdapat $\xi_1, \ldots, \xi_m \in \bigwedge^{k-1} V$ dengan $\xi = \sum_{i=1}^m v_i \wedge \xi_i$ jika dan hanya jika $v_1 \wedge \cdots \wedge v_m \wedge \xi = 0$.
	\item Buktikan bahwa operasi kontraksi dalam Definisi--Teorema \ref{def:wedge-contraction} menginduksi, untuk setiap $m \geq 0$, isomorfisme kanonik $\bigwedge^m(V^\vee) \rightiso (\bigwedge^m V)^\vee$.
\end{Exercises}
